\section{Count-based fitness calculations}
\label{app:scores}

\paragraph{PDZ.}
Let $n_{ir}^{\mathrm{in}}$ and $n_{ir}^{\mathrm{out}}$ denote counts for
variant $i$ in replicate $r$, and let $w$ index the wild-type reference. Let
$\mathcal R_i$ be the set of replicates for which all four counts for $i$
and $w$ are positive. For $r\in\mathcal R_i$, define the wild-type-relative
log enrichment $f_{ir}$ and its working variance $v_{ir}$ by
\begin{align}
f_{ir}&:=\log\frac{n_{ir}^{\mathrm{out}}}{n_{ir}^{\mathrm{in}}}
-\log\frac{n_{wr}^{\mathrm{out}}}{n_{wr}^{\mathrm{in}}},\\
v_{ir}&:=\frac{1}{n_{ir}^{\mathrm{in}}}+\frac{1}{n_{ir}^{\mathrm{out}}}
+\frac{1}{n_{wr}^{\mathrm{in}}}+\frac{1}{n_{wr}^{\mathrm{out}}},\qquad
g_i(\data):=\frac{\sum_{r\in\mathcal R_i} f_{ir}/v_{ir}}
{\sum_{r\in\mathcal R_i}1/v_{ir}}.
\end{align}
Only valid replicates enter the sums. We define the score as undefined when
$\mathcal R_i$ is empty, so undefined logarithms never enter the weighted
average. This inverse-Poisson-variance calculation is a DiMSum-style
approximation. The exact source DiMSum revision and configuration have not
been recovered, so we do not describe it as an exact reproduction of the
published score \citep{faure2020,zarin2024}.

\paragraph{TrpB.}
Following the deposited calculations of \citet{johnston2024}, we form log
output/input ratios separately for each retained timepoint and replicate.
We subtract the same-observation mean for stop-codon variants and divide by
the difference between the same-observation maximum and that stop-variant
mean. The final score averages valid normalized observations. The three-site
libraries use one input count per variant; the four-site landscape uses the
mean of two input-replicate counts. Inputs must be at least 10 and required
outputs positive; if an observation has two output replicates, both must be
positive. Library-specific timepoints follow the deposited notebooks; the
four-site calculation excludes 12 hours and retains 16, 20, 24, 28, and
36 hours. Culture-density and time factors cancel from the within-observation
stop/maximum normalization. The source count calculations reproduce the
published finite scores and coverage at numerical tolerance $10^{-12}$.

\paragraph{NucB.}
For the highest selected fluorescence gate, we calculate
\[
g_i(\data):=\frac{n_i^{\mathrm{gate}}/D^{\mathrm{gate}}}
{n_i^{\mathrm{in}}/D^{\mathrm{in}}},
\]
where $D^{\mathrm{gate}}$ and $D^{\mathrm{in}}$ are the source normalization
totals, including controls when specified by the source. The selected gates
are generation-specific \citep{thomas2025}. A zero gate count gives zero
enrichment; a zero input count makes the score undefined. Source controls
contribute to the appropriate totals but not to ranking candidates.
Generation 1 retains the source total correction for sequences omitted
from the fitted axis; this common multiplicative correction does not alter
ranks. The observed count calculations reproduce source enrichment values
at tolerance $10^{-12}$.

\section{Ranking and computational conventions}
\label{app:computation}

\paragraph{Finite samples and reference support.}
We store sampled scores as single-precision arrays and retain samples with
at least 1,000 finite candidate scores. Table~\ref{tab:inventory} gives the
resulting sample-count ranges. Rank intervals and reported-rank percentiles
use the finite observations for each candidate; undefined scores do not
become zero scores. Inclusion counts use membership in retained size-$k$
sets, with the retained sample count as denominator. Count-based median-score
references require at least two finite scores per candidate; parameter
collections require at least one. Stored sample identities preserve
alignment across parameters and count draws. Source controls and other
excluded candidates remain outside the ranking universe.

\paragraph{Ranks and selections.}
Exact tied scores receive average ranks for interval summaries. Fixed-size
top-$k$ selection uses the deterministic selection order, with source block
order retained where needed to reproduce prepared results. Because the two
conventions answer different questions at a tie, we calculate set membership
from stored selections rather than from average rank thresholds. The
median-score reference orders median scores and resolves exact ties
deterministically. Rankings always use the eligible full competitor pool;
reported cohorts and displayed summaries remain top-focused.

\paragraph{Exact shared fractions.}
Let $c_i(k):=\sum_s\mathbf{1}\{i\in T_k^{(s)}\}$ be the number of retained
sets containing variant $i$. Then sample-pair overlap is exactly
\[
\widehat C_k=\frac{\sum_i c_i(k)[c_i(k)-1]}{S(S-1)k}.
\]
We use this identity instead of drawing random pairs of samples. It is an
empirical summary of distinct retained samples; dependence among Monte
Carlo draws is not removed by the identity. We likewise obtain reference
overlap by averaging inclusion frequencies over each fixed reference set.

\subsection{What the percentile area measures}
\label{app:area-interpretation}

The area $A=\int_0^1|F(u)-u|\,du$ describes departure of the
reported-rank percentile distribution from uniformity. It is not the
ordinary area under an empirical curve, a receiver-operating-characteristic
area, or a rank error. Three simple cases distinguish its interpretation.
If every sampled rank equals its reported rank, all percentiles equal
0.5 and $A=0.25$. If half the percentiles equal zero and half equal
one, $A$ is also 0.25 despite a very different curve. If every
percentile equals zero, or every percentile equals one, $A=0.5$;
these opposite directions have the same unsigned area.

The implementation integrates each constant segment of the empirical
cumulative distribution against the diagonal, splitting a segment where
they cross. This avoids interpolation between successive jumps. Tests
verify the cases above, $A=0.125$ for percentiles $\{0.25,0.75\}$,
permutation invariance, and rejection of empty, nonfinite, or out-of-range
inputs. Each supplementary curve export recomputes the exact area from
its complete stored percentile vector and checks agreement with its
stored scalar to absolute tolerance $10^{-12}$.

Spearman correlation retains a different feature: the relative order
of reported ranks and median sampled ranks within the cohort. For
example, all candidates in a reported cohort could move down by the
same number of positions because other eligible candidates overtake
them, yet their within-cohort correlation would remain one. This is
an illustrative possibility, not a claimed result for every example.
Neither correlation nor percentile area measures the spread of each
candidate's ranks or the overlap between jointly sampled batches.

Uniformity can be a target in an appropriate repeated-data calibration
study, but the present comparison selects leading candidates and uses
reported and fitted scores derived from the same observations, sometimes
with different score definitions. The diagonal is therefore descriptive.
There is no demonstrated target value of $A$ for a correct biological model
in this analysis.

\paragraph{Implementation checks.}
Repository tests compare streamed rankings with dense calculations,
exercise ties and candidate exclusions, and check table and figure
outputs. These checks do not demonstrate frequentist interval coverage
or biological model adequacy. A controlled synthetic coverage study
is not part of the present draft.

\section{Descriptive association analyses}
\label{app:associations}

The unaggregated relationships are shown for every landscape, model,
and sampling recipe in Supplementary Section~\ref{app:per-analysis-why}.
Those scatterplots contain no fitted regression lines. The regression
coefficients summarized here are separate descriptive calculations;
the point displays let the reader assess patterns that a coefficient
or an across-landscape median can obscure.

Recall that $U_i$ is normalized score width and $\widetilde r_i$ is the
median-score reference rank. Local score spacing is the average change in
median score per rank across up to five neighbors on each side of a
reference rank, divided by $S_{\mathrm{run}}$. Endpoints use the available
neighbors. We define crowding as minus the base-10 logarithm of this
spacing, using a small numerical floor for zero spacings. Larger values
indicate closer competitors on this scale.

\paragraph{Figure 4 examples and plotted data.}
\label{app:figure-explanation}
Figure~\ref{fig:explanation} uses the parameter-resampling versions of
the same running examples as Figs.~\ref{fig:reported}--\ref{fig:selection}:
PDZ trans-1 and TrpB libB under Exponential--Exponential,
and NucB G3 under Exponential. Methods Section~\ref{methods:examples}
describes their selection. No Gamma specification enters this figure;
the complete displays below retain Gamma fits and all other analyses.

Each row includes every median-score reference top-100 member and
every outside candidate with $\hat q_i(100)\geq0.01$. The rows contain
419, 516, and 128 candidates, totaling 1,063. Every candidate appears in
all four panels of its row, without binning, point subsampling, or
pooling across analyses. Each point summarizes one candidate's samples;
it is not an individual sample. Stored input counts, output counts,
score medians, interval endpoints, and inclusion frequencies are unchanged.

Counts are transformed as $\log_{10}(n+1)$, retaining zero counts at zero.
Normalized widths use a logarithmic axis. All widths in these three
examples are positive, so no plot-only floor affects them. Count and
width limits cover all points and are common to the three rows. The
margin panels share limits within a row, but limits differ between
rows; apparent transition steepness should not be compared without
reading the axes. Inclusion frequencies use a common zero-to-one scale.

The exported variant table retains every plotted row and its coordinates.
The run index records identities, cohort sizes, sample counts, and mean
shared fractions. Inclusion and overlap come from joint rankings against
all finite eligible competitors, not from marginal intervals or rankings
restricted to the displayed cohort. Because challengers are selected
using inclusion frequencies, these associations do not describe all
variants in a landscape.

\paragraph{Observed count features.}
\label{app:count-features}
Features are fixed from the observed data and shared across sampling
recipes for a landscape. PDZ input and output support sum the three
paired replicates separately. TrpB input support uses the single input
count once for three-site libraries and the sum of the two input counts
for four-site; output support sums every retained output replicate and
model timepoint. NucB input support sums the arrays \texttt{ic1},
\texttt{ic2}, and \texttt{ic3} where present, over all observation axes.
Output support sums the low- and high-gate arrays for G1, the three
gate arrays for G2, and all four output rows for G3 and G4. Resequenced
libraries contribute their recorded reads to these descriptive totals.
None of these sums turns sequencing repeats, timepoints, or different
gates into independent biological replicates. Log input support in
regressions means $\log(1+n_i^{\mathrm{in}})$; scatterplots use
$\log_{10}(1+n)$ for display.

\paragraph{Observation discordance.}
For PDZ paired replicates, let $e_{ir}$ be the difference of smoothed
log output and input frequencies for variant $i$ and replicate $r$:
\[
e_{ir}:=\log\frac{n_{ir}^{\mathrm{out}}+\delta}
{D_r^{\mathrm{out}}+M\delta}
-\log\frac{n_{ir}^{\mathrm{in}}+\delta}{D_r^{\mathrm{in}}+M\delta},
\qquad
v_{ir}^{\mathrm{obs}}:=\frac{1}{n_{ir}^{\mathrm{in}}+\delta}
+\frac{1}{n_{ir}^{\mathrm{out}}+\delta},
\]
where $\delta=0.5$ and $D_r$ are totals on the feature-construction
count axis of length $M$. Set $w_{ir}:=1/v_{ir}^{\mathrm{obs}}$ and
$\overline e_i:=\sum_rw_{ir}e_{ir}/\sum_rw_{ir}$. Reduced heterogeneity is
\[
H_i:=\frac{1}{R-1}\sum_{r=1}^R w_{ir}(e_{ir}-\overline e_i)^2,
\]
where $R=3$ for PDZ. It measures replicate scatter relative to a working
Poisson log-count variance. This diagnostic smoothing does not change
the no-pseudocount fitness reconstructions in Supplementary
Section~\ref{app:scores}.

For TrpB observations with two output replicates, we use smoothed log
output frequencies and variance $1/(n_{i\ell r}^{\mathrm{out}}+\delta)$,
remove a separate weighted mean within each timepoint $\ell$, and sum
weighted squared residuals over timepoints. With $J$ timepoints and
$R=2$ replicates, the denominator is $J(R-1)$. This removes between-timepoint
changes before measuring within-timepoint disagreement. TrpB A--C has
no paired output contrast; NucB has no verified comparable feature in
the current analysis. Their discordance value is missing, not zero.
The regression includes it only when more than three decision-cohort
values are available, then uses complete cases.

Within the decision cohort, the score-precision regression uses
$\log(1+U_i)$ as response, with log input support and
$\log(1+\widetilde r_i)$ as predictors, adding observation discordance
where available. The rank regression uses
$\log[1+Q_{0.95}(R_i)-Q_{0.05}(R_i)]$ as response, with $U_i$,
crowding, their product, and $\log(1+\widetilde r_i)$ as predictors.
We standardize the response and predictors on complete cases before
fitting linear regressions. Coefficients are summarized within
landscapes by their median across models. Because of the interaction,
a main-effect coefficient is not a universal effect at every value
of the other predictor.

The score-uncertainty coefficient is positive in 17 of 18, 18 of 19,
and 17 of 18 landscape summaries for parameter, optimized-fit count,
and posterior count resampling, respectively. The crowding coefficient
is positive in 15 of 18, 17 of 19, and 14 of 18. The interaction is
positive in 3 of 18, 7 of 19, and 6 of 18. These are counts of fitted
coefficient signs, not uniformly positive conditional slopes. Predictors
and rank widths summarize the same samples; the regressions are not
independent tests of an effect of additional sequencing.

\subsection{Sensitivity to modeling assumptions}
\label{sec:sensitivity}

The running TrpB libB example has lower top-100 sample-pair overlap
than NucB G3 under all three recipes. Some comparisons are much
more sensitive to model choice. In PDZ trans-1, parameter-resampling
top-100 overlap is 59.80\% under Exponential--Exponential but 0.27\%
under Gamma--Exponential. The former is the running example; both
remain in the full results. The larger uncertainty
cannot be attributed to the experiment without checking the model.
Supplementary Section~\ref{app:fit-checks} records available
diagnostics and missing predictive checks.

For the batch-size comparison, the uniform-subset-adjusted shared
fraction is $(\widehat C_k-k/N)/(1-k/N)$, where $N$ is the eligible
candidate count. Applying this adjustment leaves the counts of
landscapes with greater overlap at $k=1000$ than at $k=100$
unchanged: 13/18, 14/19, and 13/18 for parameter, optimized-fit
count, and posterior count resampling. TrpB D--H remains the
common set of five decreasing landscapes. Table~\ref{tab:landscape-set-sizes}
gives the unadjusted endpoint values.

\section{Supporting figure views}

\begin{table}[p]
\centering
\caption{Mean shared top-k fractions (percent), summarized by the median over model specifications within each landscape. P, O, and PC denote parameter, optimized-fit count, and posterior-count resampling; the suffix denotes k. Missing entries indicate no analyzed posterior fit for that landscape. This is an endpoint comparison, not a claim of monotonicity over every intermediate k.}
\label{tab:landscape-set-sizes}
\begin{tabular}{lrrrrrr}
\toprule
Landscape & P100 & P1000 & O100 & O1000 & PC100 & PC1000 \\
\midrule
nuc-G1 & 68.0 & 85.3 & 88.7 & 97.3 & 91.9 & 96.4 \\
nuc-G2 & 91.7 & 97.4 & 89.3 & 97.2 & 85.7 & 96.1 \\
nuc-G3 & 93.5 & 97.8 & 87.0 & 97.5 & 87.8 & 96.5 \\
nuc-G4 & 94.7 & 98.4 & 60.3 & 97.5 & 81.6 & 97.0 \\
pdz-cis & 82.0 & 91.9 & 92.2 & 96.2 & 88.5 & 94.5 \\
pdz-cript-c & 91.8 & 96.0 & 93.8 & 95.4 & 91.4 & 94.4 \\
pdz-cript-n & 0.6 & 4.9 & 73.4 & 83.9 & 36.9 & 64.0 \\
pdz-trans-1 & 30.0 & 36.8 & 67.9 & 74.1 & 46.4 & 60.9 \\
pdz-trans-2 & 23.0 & 32.8 & 49.6 & 67.5 & 31.6 & 51.0 \\
trpb-four-site & -- & -- & 87.6 & 95.1 & -- & -- \\
trpb-libA & 51.8 & 63.1 & 28.4 & 41.7 & 24.0 & 36.3 \\
trpb-libB & 50.9 & 76.0 & 32.2 & 59.7 & 25.9 & 51.9 \\
trpb-libC & 52.2 & 73.3 & 32.5 & 55.2 & 27.7 & 48.8 \\
trpb-libD & 98.0 & 78.8 & 97.6 & 76.6 & 97.4 & 71.8 \\
trpb-libE & 66.3 & 49.5 & 60.3 & 46.2 & 54.6 & 39.0 \\
trpb-libF & 82.3 & 65.4 & 72.6 & 57.6 & 69.4 & 52.5 \\
trpb-libG & 82.3 & 59.9 & 81.1 & 60.3 & 77.7 & 53.4 \\
trpb-libH & 80.3 & 41.1 & 74.0 & 39.9 & 69.0 & 35.6 \\
trpb-libI & 64.4 & 84.7 & 74.5 & 83.6 & 70.5 & 80.8 \\
\bottomrule
\end{tabular}
\end{table}


\begin{figure}[p]
\centering
\includegraphics[width=\linewidth]{reported_order_top1000.pdf}
\caption{Reported-rank summaries for the reported top 1,000.
Panels (a)--(c) show percentile curves for PDZ trans-1 and TrpB libB
(both Exponential--Exponential), and
NucB G3 (Exponential). Panels (d)--(f) show the area between each
curve and the diagonal; panels (g)--(i) show Spearman correlation
between reported and median sampled ranks. Small filled dots summarize
individual model specifications within each landscape and sampling scheme,
not individual draws. Hollow blue circles, red squares, and green triangles
show model medians for parameter, optimized-fit count, and posterior count
resampling. Vertical offsets are for visibility only; horizontal values
are unchanged. Black outlines and arrows mark the exact example model
points. Ranks still compare each candidate with all finite
eligible competitors, not just this cohort. The diagonal is
descriptive and area is not a rank-error or calibration score.}
\label{fig:reported1000}
\end{figure}

\begin{figure}[p]
\centering
\includegraphics[width=\linewidth]{count_support.pdf}
\caption{Count support and score uncertainty across landscapes.
Panels show within-run Spearman correlations of input counts
(left) and output counts (right) with normalized central 90\%
score width. Each run uses its reference top 100 plus challengers
with inclusion frequency at least 0.01. Small filled dots each summarize
one model specification within a landscape and sampling scheme, not an
individual draw. Hollow blue circles, red squares, and green triangles
mark model medians for parameter, optimized-fit count, and posterior count
resampling. Vertical offsets improve visibility only; horizontal values
are unchanged.
These are associations in selected cohorts, not independent
experimental effects.}
\label{fig:counts}
\end{figure}

\clearpage
\section{Reported-rank curves and per-analysis statistics}
\label{app:all-reported-curves}

Table~\ref{tab:reported-curve-statistics} reports the curve area and
reported-versus-median-sampled-rank Spearman correlation for each of
the 158 analysis runs. Both reported cohorts, top 100 and top 1,000,
appear in the same row. The table also records retained samples and
supported candidate counts. IDs use P for parameter, O for optimized-fit
count, and C for posterior count resampling. Model codes use E, L,
and G for Exponential, Lomax, and Gamma priors; a second E or S
denotes exponential or sigmoid growth. NucB has only the prior code.
Rounding affects table display only; the companion CSV retains full
precision and complete run identifiers.

The 316 curves are overlaid in nine figures, one per protein system
and sampling scheme (Table~\ref{tab:reported-curves-index}). Each
figure has separate panels for the reported top 100 and top 1,000.
Color identifies the landscape, line style identifies the model, and
a thicker line marks the main example. All models, including Gamma,
are retained. Each line is the complete empirical step function for
one analysis; percentiles are neither pooled nor subsampled. Similar
curves can overlap, so the table supplies values that are difficult
to distinguish visually. The index links each run and cohort to its
figure and panel.

\begingroup
\small
\setlength{\tabcolsep}{3pt}
\begin{longtable}[c]{lllrrrrrrr}
\caption{Reported-rank summaries for every analysis. For each cohort, $n$ is the supported count, $A$ the area between curve and diagonal, and $\rho$ the Spearman correlation. IDs and model codes are defined in Section~\ref{app:all-reported-curves}; sample counts are retained draws. Statistics are rounded to three decimals; -- denotes an undefined value.} \label{tab:reported-curve-statistics} \\
\toprule
ID & Landscape & Model & Samples & $n_{100}$ & $A_{100}$ & $\rho_{100}$ & $n_{1000}$ & $A_{1000}$ & $\rho_{1000}$ \\
\midrule
\endfirsthead
\caption[]{Reported-rank summaries for every analysis. For each cohort, $n$ is the supported count, $A$ the area between curve and diagonal, and $\rho$ the Spearman correlation. IDs and model codes are defined in Section~\ref{app:all-reported-curves}; sample counts are retained draws. Statistics are rounded to three decimals; -- denotes an undefined value.} \\
\toprule
ID & Landscape & Model & Samples & $n_{100}$ & $A_{100}$ & $\rho_{100}$ & $n_{1000}$ & $A_{1000}$ & $\rho_{1000}$ \\
\midrule
\endhead
\midrule
\multicolumn{10}{r}{Continued on next page} \\
\midrule
\endfoot
\bottomrule
\endlastfoot
P01 & PDZ trans-1 & E/E & 16,000 & 100 & 0.210 & 0.310 & 1,000 & 0.128 & 0.542 \\
P02 & PDZ trans-1 & G/E & 16,000 & 100 & 0.496 & 0.378 & 1,000 & 0.475 & 0.607 \\
P03 & PDZ trans-2 & E/E & 16,000 & 100 & 0.260 & 0.395 & 1,000 & 0.141 & 0.679 \\
P04 & PDZ trans-2 & G/S & 16,000 & 100 & 0.497 & 0.454 & 1,000 & 0.482 & 0.764 \\
P05 & PDZ cis & E/E & 16,000 & 100 & 0.098 & 0.943 & 1,000 & 0.056 & 0.944 \\
P06 & PDZ cis & L/E & 16,000 & 100 & 0.107 & 0.948 & 1,000 & 0.066 & 0.951 \\
P07 & PDZ cis & E/S & 16,000 & 100 & 0.098 & 0.943 & 1,000 & 0.056 & 0.944 \\
P08 & PDZ cript-c & E/E & 16,000 & 100 & 0.078 & 0.885 & 1,000 & 0.090 & 0.924 \\
P09 & PDZ cript-c & L/E & 24,000 & 100 & 0.042 & 0.908 & 1,000 & 0.060 & 0.934 \\
P10 & PDZ cript-c & G/E & 24,000 & 100 & 0.399 & 0.887 & 1,000 & 0.267 & 0.886 \\
P11 & PDZ cript-c & E/S & 24,000 & 100 & 0.077 & 0.885 & 1,000 & 0.090 & 0.924 \\
P12 & PDZ cript-n & G/E & 24,000 & 100 & 0.493 & -0.011 & 1,000 & 0.449 & 0.580 \\
P13 & PDZ cript-n & L/S & 16,000 & 100 & 0.092 & 0.137 & 1,000 & 0.050 & 0.688 \\
P14 & PDZ cript-n & G/S & 16,000 & 100 & 0.493 & -0.012 & 1,000 & 0.449 & 0.582 \\
P15 & TrpB libA & E/E & 16,000 & 100 & 0.317 & 0.229 & 1,000 & 0.246 & 0.125 \\
P16 & TrpB libA & L/E & 16,000 & 100 & 0.312 & 0.430 & 1,000 & 0.254 & 0.286 \\
P17 & TrpB libA & E/S & 16,000 & 100 & 0.317 & 0.232 & 1,000 & 0.247 & 0.128 \\
P18 & TrpB libB & E/E & 16,000 & 100 & 0.314 & 0.056 & 1,000 & 0.157 & 0.262 \\
P19 & TrpB libB & L/E & 16,000 & 100 & 0.324 & 0.178 & 1,000 & 0.162 & 0.367 \\
P20 & TrpB libB & E/S & 16,000 & 100 & 0.314 & 0.054 & 1,000 & 0.157 & 0.264 \\
P21 & TrpB libC & E/E & 16,000 & 100 & 0.390 & -0.082 & 1,000 & 0.273 & 0.171 \\
P22 & TrpB libC & L/E & 16,000 & 100 & 0.379 & 0.088 & 1,000 & 0.272 & 0.259 \\
P23 & TrpB libC & E/S & 16,000 & 100 & 0.390 & -0.082 & 1,000 & 0.273 & 0.173 \\
P24 & TrpB libD & E/E & 16,000 & 100 & 0.195 & 0.904 & 1,000 & 0.149 & 0.862 \\
P25 & TrpB libD & L/E & 16,000 & 100 & 0.195 & 0.899 & 1,000 & 0.203 & 0.876 \\
P26 & TrpB libD & G/E & 16,000 & 100 & 0.241 & 0.899 & 1,000 & 0.297 & 0.873 \\
P27 & TrpB libE & E/E & 16,000 & 100 & 0.291 & 0.541 & 1,000 & 0.247 & 0.327 \\
P28 & TrpB libE & L/E & 16,000 & 100 & 0.283 & 0.557 & 1,000 & 0.250 & 0.436 \\
P29 & TrpB libE & G/E & 16,000 & 100 & 0.425 & 0.554 & 1,000 & 0.386 & 0.391 \\
P30 & TrpB libF & E/E & 16,000 & 100 & 0.337 & 0.448 & 1,000 & 0.272 & 0.325 \\
P31 & TrpB libF & L/E & 16,000 & 100 & 0.332 & 0.428 & 1,000 & 0.284 & 0.388 \\
P32 & TrpB libF & G/E & 16,000 & 100 & 0.397 & 0.432 & 1,000 & 0.335 & 0.362 \\
P33 & TrpB libF & E/S & 16,000 & 100 & 0.337 & 0.450 & 1,000 & 0.272 & 0.317 \\
P34 & TrpB libG & E/E & 16,000 & 100 & 0.308 & 0.763 & 1,000 & 0.285 & 0.344 \\
P35 & TrpB libG & L/E & 16,000 & 100 & 0.318 & 0.760 & 1,000 & 0.296 & 0.435 \\
P36 & TrpB libG & G/E & 16,000 & 100 & 0.344 & 0.764 & 1,000 & 0.355 & 0.398 \\
P37 & TrpB libH & E/E & 16,000 & 100 & 0.326 & 0.581 & 1,000 & 0.287 & 0.323 \\
P38 & TrpB libH & L/E & 16,000 & 100 & 0.326 & 0.594 & 1,000 & 0.276 & 0.415 \\
P39 & TrpB libH & G/E & 16,000 & 100 & 0.401 & 0.600 & 1,000 & 0.378 & 0.376 \\
P40 & TrpB libI & E/E & 16,000 & 100 & 0.139 & 0.802 & 1,000 & 0.059 & 0.945 \\
P41 & TrpB libI & L/E & 16,000 & 100 & 0.196 & 0.746 & 1,000 & 0.148 & 0.952 \\
P42 & TrpB libI & G/E & 16,000 & 100 & 0.433 & 0.753 & 1,000 & 0.271 & 0.953 \\
P43 & NucB G1 & E & 32,000 & 100 & 0.137 & 0.376 & 1,000 & 0.074 & 0.796 \\
P44 & NucB G1 & L & 32,000 & 100 & 0.106 & 0.395 & 1,000 & 0.073 & 0.797 \\
P45 & NucB G2 & E & 16,000 & 100 & 0.142 & 0.609 & 1,000 & 0.129 & 0.972 \\
P46 & NucB G2 & L & 24,000 & 100 & 0.145 & 0.594 & 1,000 & 0.128 & 0.971 \\
P47 & NucB G3 & E & 16,000 & 100 & 0.232 & 0.584 & 1,000 & 0.173 & 0.926 \\
P48 & NucB G4 & E & 16,000 & 100 & 0.344 & 0.334 & 1,000 & 0.227 & 0.867 \\
P49 & NucB G4 & L & 16,000 & 100 & 0.353 & 0.333 & 1,000 & 0.227 & 0.866 \\
O01 & PDZ cis & E/E & 1,000 & 100 & 0.073 & 0.928 & 1,000 & 0.081 & 0.943 \\
O02 & PDZ cis & L/E & 1,000 & 100 & 0.073 & 0.937 & 1,000 & 0.070 & 0.949 \\
O03 & PDZ cis & E/S & 1,000 & 100 & 0.075 & 0.852 & 1,000 & 0.139 & 0.837 \\
O04 & PDZ cript-c & E/E & 1,000 & 100 & 0.158 & 0.827 & 1,000 & 0.153 & 0.903 \\
O05 & PDZ cript-c & L/E & 1,000 & 100 & 0.147 & 0.846 & 1,000 & 0.146 & 0.909 \\
O06 & PDZ cript-c & G/E & 1,000 & 100 & 0.155 & 0.826 & 1,000 & 0.169 & 0.907 \\
O07 & PDZ cript-c & E/S & 1,000 & 100 & 0.202 & 0.246 & 1,000 & 0.236 & 0.639 \\
O08 & PDZ cript-c & L/S & 1,000 & 100 & 0.204 & 0.258 & 1,000 & 0.228 & 0.642 \\
O09 & PDZ cript-n & E/E & 1,000 & 100 & 0.197 & -0.162 & 1,000 & 0.118 & 0.384 \\
O10 & PDZ cript-n & L/E & 1,000 & 100 & 0.177 & -0.117 & 1,000 & 0.107 & 0.439 \\
O11 & PDZ cript-n & G/E & 986 & 100 & 0.203 & -0.175 & 1,000 & 0.044 & 0.383 \\
O12 & PDZ cript-n & L/S & 1,000 & 100 & 0.173 & -0.119 & 1,000 & 0.114 & 0.425 \\
O13 & PDZ cript-n & G/S & 967 & 100 & 0.218 & -0.178 & 1,000 & 0.053 & 0.371 \\
O14 & PDZ trans-1 & E/E & 1,000 & 100 & 0.290 & 0.190 & 1,000 & 0.213 & 0.377 \\
O15 & PDZ trans-1 & G/E & 1,000 & 100 & 0.271 & 0.213 & 1,000 & 0.190 & 0.400 \\
O16 & PDZ trans-2 & E/E & 1,000 & 100 & 0.318 & 0.255 & 1,000 & 0.204 & 0.544 \\
O17 & PDZ trans-2 & G/E & 1,000 & 100 & 0.299 & 0.290 & 1,000 & 0.178 & 0.580 \\
O18 & PDZ trans-2 & L/S & 1,000 & 100 & 0.298 & 0.284 & 1,000 & 0.176 & 0.583 \\
O19 & PDZ trans-2 & G/S & 1,000 & 100 & 0.300 & 0.296 & 1,000 & 0.178 & 0.579 \\
O20 & TrpB libA & E/E & 1,000 & 100 & 0.365 & 0.209 & 1,000 & 0.294 & 0.104 \\
O21 & TrpB libA & L/E & 1,000 & 100 & 0.348 & 0.418 & 1,000 & 0.287 & 0.276 \\
O22 & TrpB libA & E/S & 1,000 & 100 & 0.363 & 0.187 & 1,000 & 0.294 & 0.105 \\
O23 & TrpB libB & E/E & 1,000 & 100 & 0.349 & 0.052 & 1,000 & 0.191 & 0.252 \\
O24 & TrpB libB & L/E & 1,000 & 100 & 0.351 & 0.169 & 1,000 & 0.198 & 0.356 \\
O25 & TrpB libB & E/S & 1,000 & 100 & 0.348 & 0.042 & 1,000 & 0.191 & 0.252 \\
O26 & TrpB libC & E/E & 1,000 & 100 & 0.402 & -0.089 & 1,000 & 0.293 & 0.147 \\
O27 & TrpB libC & L/E & 1,000 & 100 & 0.389 & 0.069 & 1,000 & 0.290 & 0.250 \\
O28 & TrpB libC & E/S & 1,000 & 100 & 0.404 & -0.093 & 1,000 & 0.293 & 0.152 \\
O29 & TrpB libD & E/E & 1,000 & 100 & 0.198 & 0.902 & 1,000 & 0.182 & 0.865 \\
O30 & TrpB libD & L/E & 1,000 & 100 & 0.198 & 0.900 & 1,000 & 0.216 & 0.875 \\
O31 & TrpB libD & G/E & 1,000 & 100 & 0.197 & 0.900 & 1,000 & 0.195 & 0.873 \\
O32 & TrpB libE & E/E & 1,000 & 100 & 0.290 & 0.541 & 1,000 & 0.267 & 0.347 \\
O33 & TrpB libE & L/E & 1,000 & 100 & 0.272 & 0.555 & 1,000 & 0.254 & 0.454 \\
O34 & TrpB libE & G/E & 1,000 & 100 & 0.283 & 0.545 & 1,000 & 0.261 & 0.356 \\
O35 & TrpB libF & E/E & 1,000 & 99 & 0.320 & 0.440 & 999 & 0.265 & 0.319 \\
O36 & TrpB libF & L/E & 1,000 & 99 & 0.312 & 0.429 & 999 & 0.249 & 0.399 \\
O37 & TrpB libF & G/E & 1,000 & 99 & 0.313 & 0.434 & 999 & 0.258 & 0.312 \\
O38 & TrpB libF & E/S & 1,000 & 99 & 0.323 & 0.477 & 999 & 0.264 & 0.313 \\
O39 & TrpB libG & E/E & 1,000 & 100 & 0.289 & 0.767 & 1,000 & 0.283 & 0.330 \\
O40 & TrpB libG & L/E & 1,000 & 99 & 0.288 & 0.775 & 999 & 0.263 & 0.447 \\
O41 & TrpB libG & G/E & 1,000 & 100 & 0.281 & 0.772 & 1,000 & 0.273 & 0.332 \\
O42 & TrpB libH & E/E & 1,000 & 100 & 0.323 & 0.577 & 1,000 & 0.301 & 0.325 \\
O43 & TrpB libH & L/E & 1,000 & 100 & 0.316 & 0.598 & 1,000 & 0.279 & 0.436 \\
O44 & TrpB libH & G/E & 1,000 & 100 & 0.318 & 0.583 & 1,000 & 0.302 & 0.303 \\
O45 & TrpB libI & E/E & 1,000 & 100 & 0.174 & 0.800 & 1,000 & 0.101 & 0.947 \\
O46 & TrpB libI & L/E & 1,000 & 100 & 0.234 & 0.739 & 1,000 & 0.176 & 0.956 \\
O47 & TrpB libI & G/E & 1,000 & 100 & 0.222 & 0.740 & 1,000 & 0.143 & 0.955 \\
O48 & TrpB four-site & E/E & 1,000 & 100 & 0.210 & 0.770 & 1,000 & 0.109 & 0.928 \\
O49 & TrpB four-site & L/E & 1,000 & 100 & 0.210 & 0.774 & 1,000 & 0.108 & 0.927 \\
O50 & TrpB four-site & G/E & 1,000 & 100 & 0.209 & 0.770 & 1,000 & 0.108 & 0.927 \\
O51 & TrpB four-site & E/S & 1,000 & 100 & 0.202 & 0.633 & 1,000 & 0.155 & 0.743 \\
O52 & TrpB four-site & L/S & 1,000 & 100 & 0.202 & 0.633 & 1,000 & 0.152 & 0.743 \\
O53 & TrpB four-site & G/S & 1,000 & 99 & 0.145 & 0.829 & 995 & 0.213 & 0.971 \\
O54 & NucB G1 & E & 1,000 & 100 & 0.136 & 0.762 & 1,000 & 0.095 & 0.979 \\
O55 & NucB G1 & L & 1,000 & 100 & 0.158 & 0.724 & 1,000 & 0.095 & 0.975 \\
O56 & NucB G2 & E & 1,000 & 100 & 0.129 & 0.592 & 1,000 & 0.093 & 0.979 \\
O57 & NucB G2 & L & 1,000 & 100 & 0.139 & 0.557 & 1,000 & 0.095 & 0.978 \\
O58 & NucB G3 & E & 1,000 & 100 & 0.229 & 0.503 & 1,000 & 0.154 & 0.921 \\
O59 & NucB G4 & E & 1,000 & 100 & 0.362 & 0.308 & 1,000 & 0.211 & 0.867 \\
O60 & NucB G4 & L & 1,000 & 100 & 0.410 & 0.116 & 1,000 & 0.192 & 0.823 \\
C01 & PDZ trans-1 & E/E & 16,000 & 100 & 0.237 & 0.319 & 1,000 & 0.153 & 0.544 \\
C02 & PDZ trans-1 & G/E & 16,000 & 100 & 0.240 & 0.382 & 1,000 & 0.153 & 0.621 \\
C03 & PDZ trans-2 & E/E & 16,000 & 100 & 0.299 & 0.396 & 1,000 & 0.183 & 0.676 \\
C04 & PDZ trans-2 & G/S & 16,000 & 100 & 0.302 & 0.471 & 1,000 & 0.207 & 0.763 \\
C05 & PDZ cis & E/E & 16,000 & 100 & 0.081 & 0.950 & 1,000 & 0.038 & 0.956 \\
C06 & PDZ cis & L/E & 16,000 & 100 & 0.104 & 0.956 & 1,000 & 0.061 & 0.958 \\
C07 & PDZ cis & E/S & 16,000 & 100 & 0.081 & 0.950 & 1,000 & 0.038 & 0.956 \\
C08 & PDZ cript-c & E/E & 16,000 & 100 & 0.064 & 0.883 & 1,000 & 0.112 & 0.920 \\
C09 & PDZ cript-c & L/E & 24,000 & 100 & 0.030 & 0.903 & 1,000 & 0.135 & 0.927 \\
C10 & PDZ cript-c & G/E & 24,000 & 100 & 0.042 & 0.893 & 1,000 & 0.081 & 0.937 \\
C11 & PDZ cript-c & E/S & 24,000 & 100 & 0.064 & 0.883 & 1,000 & 0.112 & 0.920 \\
C12 & PDZ cript-n & G/E & 23,685 & 100 & 0.254 & -0.018 & 1,000 & 0.119 & 0.545 \\
C13 & PDZ cript-n & L/S & 15,918 & 100 & 0.145 & 0.133 & 1,000 & 0.071 & 0.682 \\
C14 & PDZ cript-n & G/S & 15,839 & 100 & 0.253 & -0.017 & 1,000 & 0.119 & 0.545 \\
C15 & TrpB libA & E/E & 16,000 & 100 & 0.378 & 0.228 & 1,000 & 0.308 & 0.133 \\
C16 & TrpB libA & L/E & 16,000 & 100 & 0.371 & 0.413 & 1,000 & 0.304 & 0.282 \\
C17 & TrpB libA & E/S & 16,000 & 100 & 0.378 & 0.230 & 1,000 & 0.308 & 0.135 \\
C18 & TrpB libB & E/E & 16,000 & 100 & 0.364 & 0.050 & 1,000 & 0.216 & 0.262 \\
C19 & TrpB libB & L/E & 16,000 & 100 & 0.372 & 0.169 & 1,000 & 0.225 & 0.363 \\
C20 & TrpB libB & E/S & 16,000 & 100 & 0.364 & 0.054 & 1,000 & 0.216 & 0.263 \\
C21 & TrpB libC & E/E & 16,000 & 100 & 0.406 & -0.079 & 1,000 & 0.302 & 0.175 \\
C22 & TrpB libC & L/E & 16,000 & 100 & 0.396 & 0.080 & 1,000 & 0.300 & 0.258 \\
C23 & TrpB libC & E/S & 16,000 & 100 & 0.406 & -0.073 & 1,000 & 0.302 & 0.176 \\
C24 & TrpB libD & E/E & 16,000 & 100 & 0.196 & 0.904 & 1,000 & 0.192 & 0.860 \\
C25 & TrpB libD & L/E & 16,000 & 100 & 0.196 & 0.900 & 1,000 & 0.223 & 0.876 \\
C26 & TrpB libD & G/E & 16,000 & 100 & 0.196 & 0.900 & 1,000 & 0.210 & 0.873 \\
C27 & TrpB libE & E/E & 16,000 & 100 & 0.291 & 0.543 & 1,000 & 0.287 & 0.329 \\
C28 & TrpB libE & L/E & 16,000 & 100 & 0.276 & 0.558 & 1,000 & 0.271 & 0.432 \\
C29 & TrpB libE & G/E & 16,000 & 100 & 0.282 & 0.553 & 1,000 & 0.278 & 0.389 \\
C30 & TrpB libF & E/E & 16,000 & 100 & 0.316 & 0.458 & 1,000 & 0.268 & 0.317 \\
C31 & TrpB libF & L/E & 16,000 & 99 & 0.312 & 0.433 & 999 & 0.251 & 0.376 \\
C32 & TrpB libF & G/E & 16,000 & 100 & 0.307 & 0.448 & 1,000 & 0.257 & 0.354 \\
C33 & TrpB libF & E/S & 16,000 & 100 & 0.313 & 0.459 & 1,000 & 0.268 & 0.313 \\
C34 & TrpB libG & E/E & 16,000 & 100 & 0.288 & 0.769 & 1,000 & 0.290 & 0.323 \\
C35 & TrpB libG & L/E & 16,000 & 100 & 0.283 & 0.771 & 1,000 & 0.269 & 0.417 \\
C36 & TrpB libG & G/E & 16,000 & 100 & 0.283 & 0.775 & 1,000 & 0.277 & 0.380 \\
C37 & TrpB libH & E/E & 16,000 & 100 & 0.323 & 0.579 & 1,000 & 0.316 & 0.314 \\
C38 & TrpB libH & L/E & 16,000 & 100 & 0.319 & 0.594 & 1,000 & 0.294 & 0.410 \\
C39 & TrpB libH & G/E & 16,000 & 100 & 0.319 & 0.598 & 1,000 & 0.308 & 0.370 \\
C40 & TrpB libI & E/E & 16,000 & 100 & 0.193 & 0.804 & 1,000 & 0.114 & 0.944 \\
C41 & TrpB libI & L/E & 16,000 & 100 & 0.255 & 0.741 & 1,000 & 0.187 & 0.956 \\
C42 & TrpB libI & G/E & 16,000 & 100 & 0.240 & 0.748 & 1,000 & 0.161 & 0.955 \\
C43 & NucB G1 & E & 32,000 & 100 & 0.041 & 0.909 & 1,000 & 0.048 & 0.992 \\
C44 & NucB G1 & L & 32,000 & 100 & 0.027 & 0.916 & 1,000 & 0.051 & 0.992 \\
C45 & NucB G2 & E & 16,000 & 100 & 0.151 & 0.560 & 1,000 & 0.060 & 0.977 \\
C46 & NucB G2 & L & 24,000 & 100 & 0.172 & 0.475 & 1,000 & 0.065 & 0.976 \\
C47 & NucB G3 & E & 16,000 & 100 & 0.235 & 0.564 & 1,000 & 0.136 & 0.926 \\
C48 & NucB G4 & E & 16,000 & 100 & 0.363 & 0.326 & 1,000 & 0.204 & 0.865 \\
C49 & NucB G4 & L & 16,000 & 100 & 0.365 & 0.323 & 1,000 & 0.204 & 0.863 \\
\end{longtable}
\endgroup

\clearpage

\begin{table}[htbp]
\centering
\small
\caption{Grouped reported-rank percentile curves. Each figure overlays one protein system and sampling scheme, with reported cohorts in separate panels.}
\label{tab:reported-curves-index}
\begin{tabular}{llrr}
\toprule
Scheme & Family & Analyses & Figures \\
\midrule
Parameter & PDZ & 14 & \ref{fig:curves-parameter_posterior-pdz-1} \\
Parameter & TrpB & 28 & \ref{fig:curves-parameter_posterior-trpb-1} \\
Parameter & NucB & 7 & \ref{fig:curves-parameter_posterior-nuc-1} \\
Optimized-fit counts & PDZ & 19 & \ref{fig:curves-mle_predictive-pdz-1} \\
Optimized-fit counts & TrpB & 34 & \ref{fig:curves-mle_predictive-trpb-1} \\
Optimized-fit counts & NucB & 7 & \ref{fig:curves-mle_predictive-nuc-1} \\
Posterior counts & PDZ & 14 & \ref{fig:curves-posterior_predictive-pdz-1} \\
Posterior counts & TrpB & 28 & \ref{fig:curves-posterior_predictive-trpb-1} \\
Posterior counts & NucB & 7 & \ref{fig:curves-posterior_predictive-nuc-1} \\
\bottomrule
\end{tabular}
\end{table}
\clearpage

\begin{figure}[p]
\centering
\includegraphics[width=\linewidth]{reported_curves_parameter_posterior_pdz_01.pdf}
\caption{PDZ reported-rank percentile curves: parameter. Each of the 14 lines per panel is one landscape/model analysis; panels use the reported top 100 and top 1000. Colors identify landscapes and line styles identify models. E, L, and G denote Exponential, Lomax, and Gamma priors; a second E or S denotes exponential or sigmoid growth. A thicker line marks the main example. All percentiles are retained; curves are overlaid, not averaged or pooled. Table~\ref{tab:reported-curve-statistics} gives each analysis's area, correlation, and counts. The gray diagonal is descriptive, not a calibration target.}
\label{fig:curves-parameter_posterior-pdz-1}
\end{figure}
\clearpage

\begin{figure}[p]
\centering
\includegraphics[width=\linewidth]{reported_curves_parameter_posterior_trpb_01.pdf}
\caption{TrpB reported-rank percentile curves: parameter. Each of the 28 lines per panel is one landscape/model analysis; panels use the reported top 100 and top 1000. Colors identify landscapes and line styles identify models. E, L, and G denote Exponential, Lomax, and Gamma priors; a second E or S denotes exponential or sigmoid growth. A thicker line marks the main example. All percentiles are retained; curves are overlaid, not averaged or pooled. Table~\ref{tab:reported-curve-statistics} gives each analysis's area, correlation, and counts. The gray diagonal is descriptive, not a calibration target.}
\label{fig:curves-parameter_posterior-trpb-1}
\end{figure}
\clearpage

\begin{figure}[p]
\centering
\includegraphics[width=\linewidth]{reported_curves_parameter_posterior_nuc_01.pdf}
\caption{NucB reported-rank percentile curves: parameter. Each of the 7 lines per panel is one landscape/model analysis; panels use the reported top 100 and top 1000. Colors identify landscapes and line styles identify models. E, L, and G denote Exponential, Lomax, and Gamma priors; a second E or S denotes exponential or sigmoid growth. A thicker line marks the main example. All percentiles are retained; curves are overlaid, not averaged or pooled. Table~\ref{tab:reported-curve-statistics} gives each analysis's area, correlation, and counts. The gray diagonal is descriptive, not a calibration target.}
\label{fig:curves-parameter_posterior-nuc-1}
\end{figure}
\clearpage

\begin{figure}[p]
\centering
\includegraphics[width=\linewidth]{reported_curves_mle_predictive_pdz_01.pdf}
\caption{PDZ reported-rank percentile curves: optimized-fit counts. Each of the 19 lines per panel is one landscape/model analysis; panels use the reported top 100 and top 1000. Colors identify landscapes and line styles identify models. E, L, and G denote Exponential, Lomax, and Gamma priors; a second E or S denotes exponential or sigmoid growth. A thicker line marks the main example. All percentiles are retained; curves are overlaid, not averaged or pooled. Table~\ref{tab:reported-curve-statistics} gives each analysis's area, correlation, and counts. The gray diagonal is descriptive, not a calibration target.}
\label{fig:curves-mle_predictive-pdz-1}
\end{figure}
\clearpage

\begin{figure}[p]
\centering
\includegraphics[width=\linewidth]{reported_curves_mle_predictive_trpb_01.pdf}
\caption{TrpB reported-rank percentile curves: optimized-fit counts. Each of the 34 lines per panel is one landscape/model analysis; panels use the reported top 100 and top 1000. Colors identify landscapes and line styles identify models. E, L, and G denote Exponential, Lomax, and Gamma priors; a second E or S denotes exponential or sigmoid growth. A thicker line marks the main example. All percentiles are retained; curves are overlaid, not averaged or pooled. Table~\ref{tab:reported-curve-statistics} gives each analysis's area, correlation, and counts. The gray diagonal is descriptive, not a calibration target.}
\label{fig:curves-mle_predictive-trpb-1}
\end{figure}
\clearpage

\begin{figure}[p]
\centering
\includegraphics[width=\linewidth]{reported_curves_mle_predictive_nuc_01.pdf}
\caption{NucB reported-rank percentile curves: optimized-fit counts. Each of the 7 lines per panel is one landscape/model analysis; panels use the reported top 100 and top 1000. Colors identify landscapes and line styles identify models. E, L, and G denote Exponential, Lomax, and Gamma priors; a second E or S denotes exponential or sigmoid growth. A thicker line marks the main example. All percentiles are retained; curves are overlaid, not averaged or pooled. Table~\ref{tab:reported-curve-statistics} gives each analysis's area, correlation, and counts. The gray diagonal is descriptive, not a calibration target.}
\label{fig:curves-mle_predictive-nuc-1}
\end{figure}
\clearpage

\begin{figure}[p]
\centering
\includegraphics[width=\linewidth]{reported_curves_posterior_predictive_pdz_01.pdf}
\caption{PDZ reported-rank percentile curves: posterior counts. Each of the 14 lines per panel is one landscape/model analysis; panels use the reported top 100 and top 1000. Colors identify landscapes and line styles identify models. E, L, and G denote Exponential, Lomax, and Gamma priors; a second E or S denotes exponential or sigmoid growth. A thicker line marks the main example. All percentiles are retained; curves are overlaid, not averaged or pooled. Table~\ref{tab:reported-curve-statistics} gives each analysis's area, correlation, and counts. The gray diagonal is descriptive, not a calibration target.}
\label{fig:curves-posterior_predictive-pdz-1}
\end{figure}
\clearpage

\begin{figure}[p]
\centering
\includegraphics[width=\linewidth]{reported_curves_posterior_predictive_trpb_01.pdf}
\caption{TrpB reported-rank percentile curves: posterior counts. Each of the 28 lines per panel is one landscape/model analysis; panels use the reported top 100 and top 1000. Colors identify landscapes and line styles identify models. E, L, and G denote Exponential, Lomax, and Gamma priors; a second E or S denotes exponential or sigmoid growth. A thicker line marks the main example. All percentiles are retained; curves are overlaid, not averaged or pooled. Table~\ref{tab:reported-curve-statistics} gives each analysis's area, correlation, and counts. The gray diagonal is descriptive, not a calibration target.}
\label{fig:curves-posterior_predictive-trpb-1}
\end{figure}
\clearpage

\begin{figure}[p]
\centering
\includegraphics[width=\linewidth]{reported_curves_posterior_predictive_nuc_01.pdf}
\caption{NucB reported-rank percentile curves: posterior counts. Each of the 7 lines per panel is one landscape/model analysis; panels use the reported top 100 and top 1000. Colors identify landscapes and line styles identify models. E, L, and G denote Exponential, Lomax, and Gamma priors; a second E or S denotes exponential or sigmoid growth. A thicker line marks the main example. All percentiles are retained; curves are overlaid, not averaged or pooled. Table~\ref{tab:reported-curve-statistics} gives each analysis's area, correlation, and counts. The gray diagonal is descriptive, not a calibration target.}
\label{fig:curves-posterior_predictive-nuc-1}
\end{figure}
\clearpage


\clearpage
\section{Per-analysis count support and selection uncertainty}
\label{app:per-analysis-why}

We show each of the 158 runs separately, grouped by protein system
and sampling recipe. Table~\ref{tab:why-index} locates the 58 figure
pages. Row headings give the landscape, model, retained sample
count, decision-cohort size, and mean shared top-100 fraction.
Each row contains the reference top 100 and all outside candidates
with inclusion frequency at least 0.01. These are the top-focused
cohorts used in Section~\ref{sec:mechanisms}, not samples of all
eligible variants.

From left to right, panels show normalized score width against
input support, output support, and score margin from the top-100
cutoff, followed by top-100 inclusion against margin. Score width
and margin are divided by the run's top-1,000 median-score IQR.
Counts use $\log_{10}(1+n)$, and score widths use logarithmic axes.
Count and width limits are common within each protein-system
and recipe collection, including across pages. Margin panels
share a scale within a row but can differ between runs.

All points are retained. Exact-zero score widths are displayed
at half the smallest positive normalized width in the corresponding
protein-system and recipe collection; numerical tables retain
zero. The index records these plot-only floors and affected point
counts. Blue filled points denote reference members and orange
open points denote challengers. Dotted lines mark zero margin.
No regression lines or smoothing curves are fitted. The figures
show within-analysis associations, not causal effects of read
depth or model-adequacy checks against new observations.

\begin{table}[htbp]
\centering
\small
\caption{Index of per-analysis count-support and selection summaries. Every one of the 158 analysis runs appears on exactly one page.}
\label{tab:why-index}
\begin{tabular}{llrr}
\toprule
Scheme & Family & Runs & Figures \\
\midrule
Parameter resampling & PDZ & 14 & \ref{fig:why-parameter_posterior-pdz-1}--\ref{fig:why-parameter_posterior-pdz-5} \\
Parameter resampling & TrpB & 28 & \ref{fig:why-parameter_posterior-trpb-1}--\ref{fig:why-parameter_posterior-trpb-10} \\
Parameter resampling & NucB & 7 & \ref{fig:why-parameter_posterior-nuc-1}--\ref{fig:why-parameter_posterior-nuc-3} \\
Optimized-fit count resampling & PDZ & 19 & \ref{fig:why-mle_predictive-pdz-1}--\ref{fig:why-mle_predictive-pdz-7} \\
Optimized-fit count resampling & TrpB & 34 & \ref{fig:why-mle_predictive-trpb-1}--\ref{fig:why-mle_predictive-trpb-12} \\
Optimized-fit count resampling & NucB & 7 & \ref{fig:why-mle_predictive-nuc-1}--\ref{fig:why-mle_predictive-nuc-3} \\
Posterior count resampling & PDZ & 14 & \ref{fig:why-posterior_predictive-pdz-1}--\ref{fig:why-posterior_predictive-pdz-5} \\
Posterior count resampling & TrpB & 28 & \ref{fig:why-posterior_predictive-trpb-1}--\ref{fig:why-posterior_predictive-trpb-10} \\
Posterior count resampling & NucB & 7 & \ref{fig:why-posterior_predictive-nuc-1}--\ref{fig:why-posterior_predictive-nuc-3} \\
\bottomrule
\end{tabular}
\end{table}
\clearpage

\begin{figure}[p]
\centering
\includegraphics[width=\linewidth]{summary_why_parameter_posterior_pdz_01.pdf}
\caption{PDZ decision-cohort summaries: parameter resampling. Rows, top to bottom: pdz-trans-1, Exponential--Exponential; pdz-trans-1, Gamma--Exponential; pdz-trans-2, Exponential--Exponential. Columns show normalized central 90\% score width versus input count, output count, and median-score margin from the top-100 cutoff; then top-100 inclusion versus that margin. Widths and margins use the run's top-1,000 median-score interquartile range (IQR). Blue filled points are reference top-100 members; orange open points are outside candidates with inclusion probability at least 0.01. Count and score-width limits are shared within family and scheme; margin limits are shared within a row. Dotted lines mark the cutoff. No regression lines are fitted.}
\label{fig:why-parameter_posterior-pdz-1}
\end{figure}
\clearpage

\begin{figure}[p]
\centering
\includegraphics[width=\linewidth]{summary_why_parameter_posterior_pdz_02.pdf}
\caption{PDZ decision-cohort summaries: parameter resampling. Rows, top to bottom: pdz-trans-2, Gamma--Sigmoid; pdz-cis, Exponential--Exponential; pdz-cis, Lomax--Exponential. Columns show normalized central 90\% score width versus input count, output count, and median-score margin from the top-100 cutoff; then top-100 inclusion versus that margin. Widths and margins use the run's top-1,000 median-score interquartile range (IQR). Blue filled points are reference top-100 members; orange open points are outside candidates with inclusion probability at least 0.01. Count and score-width limits are shared within family and scheme; margin limits are shared within a row. Dotted lines mark the cutoff. No regression lines are fitted.}
\label{fig:why-parameter_posterior-pdz-2}
\end{figure}
\clearpage

\begin{figure}[p]
\centering
\includegraphics[width=\linewidth]{summary_why_parameter_posterior_pdz_03.pdf}
\caption{PDZ decision-cohort summaries: parameter resampling. Rows, top to bottom: pdz-cis, Exponential--Sigmoid; pdz-cript-c, Exponential--Exponential; pdz-cript-c, Lomax--Exponential. Columns show normalized central 90\% score width versus input count, output count, and median-score margin from the top-100 cutoff; then top-100 inclusion versus that margin. Widths and margins use the run's top-1,000 median-score interquartile range (IQR). Blue filled points are reference top-100 members; orange open points are outside candidates with inclusion probability at least 0.01. Count and score-width limits are shared within family and scheme; margin limits are shared within a row. Dotted lines mark the cutoff. No regression lines are fitted.}
\label{fig:why-parameter_posterior-pdz-3}
\end{figure}
\clearpage

\begin{figure}[p]
\centering
\includegraphics[width=\linewidth]{summary_why_parameter_posterior_pdz_04.pdf}
\caption{PDZ decision-cohort summaries: parameter resampling. Rows, top to bottom: pdz-cript-c, Gamma--Exponential; pdz-cript-c, Exponential--Sigmoid; pdz-cript-n, Gamma--Exponential. Columns show normalized central 90\% score width versus input count, output count, and median-score margin from the top-100 cutoff; then top-100 inclusion versus that margin. Widths and margins use the run's top-1,000 median-score interquartile range (IQR). Blue filled points are reference top-100 members; orange open points are outside candidates with inclusion probability at least 0.01. Count and score-width limits are shared within family and scheme; margin limits are shared within a row. Dotted lines mark the cutoff. No regression lines are fitted.}
\label{fig:why-parameter_posterior-pdz-4}
\end{figure}
\clearpage

\begin{figure}[p]
\centering
\includegraphics[width=\linewidth]{summary_why_parameter_posterior_pdz_05.pdf}
\caption{PDZ decision-cohort summaries: parameter resampling. Rows, top to bottom: pdz-cript-n, Lomax--Sigmoid; pdz-cript-n, Gamma--Sigmoid. Columns show normalized central 90\% score width versus input count, output count, and median-score margin from the top-100 cutoff; then top-100 inclusion versus that margin. Widths and margins use the run's top-1,000 median-score interquartile range (IQR). Blue filled points are reference top-100 members; orange open points are outside candidates with inclusion probability at least 0.01. Count and score-width limits are shared within family and scheme; margin limits are shared within a row. Dotted lines mark the cutoff. No regression lines are fitted.}
\label{fig:why-parameter_posterior-pdz-5}
\end{figure}
\clearpage

\begin{figure}[p]
\centering
\includegraphics[width=\linewidth]{summary_why_parameter_posterior_trpb_01.pdf}
\caption{TrpB decision-cohort summaries: parameter resampling. Rows, top to bottom: trpb-libA, Exponential--Exponential; trpb-libA, Lomax--Exponential; trpb-libA, Exponential--Sigmoid. Columns show normalized central 90\% score width versus input count, output count, and median-score margin from the top-100 cutoff; then top-100 inclusion versus that margin. Widths and margins use the run's top-1,000 median-score interquartile range (IQR). Blue filled points are reference top-100 members; orange open points are outside candidates with inclusion probability at least 0.01. Count and score-width limits are shared within family and scheme; margin limits are shared within a row. Dotted lines mark the cutoff. No regression lines are fitted.}
\label{fig:why-parameter_posterior-trpb-1}
\end{figure}
\clearpage

\begin{figure}[p]
\centering
\includegraphics[width=\linewidth]{summary_why_parameter_posterior_trpb_02.pdf}
\caption{TrpB decision-cohort summaries: parameter resampling. Rows, top to bottom: trpb-libB, Exponential--Exponential; trpb-libB, Lomax--Exponential; trpb-libB, Exponential--Sigmoid. Columns show normalized central 90\% score width versus input count, output count, and median-score margin from the top-100 cutoff; then top-100 inclusion versus that margin. Widths and margins use the run's top-1,000 median-score interquartile range (IQR). Blue filled points are reference top-100 members; orange open points are outside candidates with inclusion probability at least 0.01. Count and score-width limits are shared within family and scheme; margin limits are shared within a row. Dotted lines mark the cutoff. No regression lines are fitted.}
\label{fig:why-parameter_posterior-trpb-2}
\end{figure}
\clearpage

\begin{figure}[p]
\centering
\includegraphics[width=\linewidth]{summary_why_parameter_posterior_trpb_03.pdf}
\caption{TrpB decision-cohort summaries: parameter resampling. Rows, top to bottom: trpb-libC, Exponential--Exponential; trpb-libC, Lomax--Exponential; trpb-libC, Exponential--Sigmoid. Columns show normalized central 90\% score width versus input count, output count, and median-score margin from the top-100 cutoff; then top-100 inclusion versus that margin. Widths and margins use the run's top-1,000 median-score interquartile range (IQR). Blue filled points are reference top-100 members; orange open points are outside candidates with inclusion probability at least 0.01. Count and score-width limits are shared within family and scheme; margin limits are shared within a row. Dotted lines mark the cutoff. No regression lines are fitted.}
\label{fig:why-parameter_posterior-trpb-3}
\end{figure}
\clearpage

\begin{figure}[p]
\centering
\includegraphics[width=\linewidth]{summary_why_parameter_posterior_trpb_04.pdf}
\caption{TrpB decision-cohort summaries: parameter resampling. Rows, top to bottom: trpb-libD, Exponential--Exponential; trpb-libD, Lomax--Exponential; trpb-libD, Gamma--Exponential. Columns show normalized central 90\% score width versus input count, output count, and median-score margin from the top-100 cutoff; then top-100 inclusion versus that margin. Widths and margins use the run's top-1,000 median-score interquartile range (IQR). Blue filled points are reference top-100 members; orange open points are outside candidates with inclusion probability at least 0.01. Count and score-width limits are shared within family and scheme; margin limits are shared within a row. Dotted lines mark the cutoff. No regression lines are fitted.}
\label{fig:why-parameter_posterior-trpb-4}
\end{figure}
\clearpage

\begin{figure}[p]
\centering
\includegraphics[width=\linewidth]{summary_why_parameter_posterior_trpb_05.pdf}
\caption{TrpB decision-cohort summaries: parameter resampling. Rows, top to bottom: trpb-libE, Exponential--Exponential; trpb-libE, Lomax--Exponential; trpb-libE, Gamma--Exponential. Columns show normalized central 90\% score width versus input count, output count, and median-score margin from the top-100 cutoff; then top-100 inclusion versus that margin. Widths and margins use the run's top-1,000 median-score interquartile range (IQR). Blue filled points are reference top-100 members; orange open points are outside candidates with inclusion probability at least 0.01. Count and score-width limits are shared within family and scheme; margin limits are shared within a row. Dotted lines mark the cutoff. No regression lines are fitted.}
\label{fig:why-parameter_posterior-trpb-5}
\end{figure}
\clearpage

\begin{figure}[p]
\centering
\includegraphics[width=\linewidth]{summary_why_parameter_posterior_trpb_06.pdf}
\caption{TrpB decision-cohort summaries: parameter resampling. Rows, top to bottom: trpb-libF, Exponential--Exponential; trpb-libF, Lomax--Exponential; trpb-libF, Gamma--Exponential. Columns show normalized central 90\% score width versus input count, output count, and median-score margin from the top-100 cutoff; then top-100 inclusion versus that margin. Widths and margins use the run's top-1,000 median-score interquartile range (IQR). Blue filled points are reference top-100 members; orange open points are outside candidates with inclusion probability at least 0.01. Count and score-width limits are shared within family and scheme; margin limits are shared within a row. Dotted lines mark the cutoff. No regression lines are fitted.}
\label{fig:why-parameter_posterior-trpb-6}
\end{figure}
\clearpage

\begin{figure}[p]
\centering
\includegraphics[width=\linewidth]{summary_why_parameter_posterior_trpb_07.pdf}
\caption{TrpB decision-cohort summaries: parameter resampling. Rows, top to bottom: trpb-libF, Exponential--Sigmoid; trpb-libG, Exponential--Exponential; trpb-libG, Lomax--Exponential. Columns show normalized central 90\% score width versus input count, output count, and median-score margin from the top-100 cutoff; then top-100 inclusion versus that margin. Widths and margins use the run's top-1,000 median-score interquartile range (IQR). Blue filled points are reference top-100 members; orange open points are outside candidates with inclusion probability at least 0.01. Count and score-width limits are shared within family and scheme; margin limits are shared within a row. Dotted lines mark the cutoff. No regression lines are fitted.}
\label{fig:why-parameter_posterior-trpb-7}
\end{figure}
\clearpage

\begin{figure}[p]
\centering
\includegraphics[width=\linewidth]{summary_why_parameter_posterior_trpb_08.pdf}
\caption{TrpB decision-cohort summaries: parameter resampling. Rows, top to bottom: trpb-libG, Gamma--Exponential; trpb-libH, Exponential--Exponential; trpb-libH, Lomax--Exponential. Columns show normalized central 90\% score width versus input count, output count, and median-score margin from the top-100 cutoff; then top-100 inclusion versus that margin. Widths and margins use the run's top-1,000 median-score interquartile range (IQR). Blue filled points are reference top-100 members; orange open points are outside candidates with inclusion probability at least 0.01. Count and score-width limits are shared within family and scheme; margin limits are shared within a row. Dotted lines mark the cutoff. No regression lines are fitted.}
\label{fig:why-parameter_posterior-trpb-8}
\end{figure}
\clearpage

\begin{figure}[p]
\centering
\includegraphics[width=\linewidth]{summary_why_parameter_posterior_trpb_09.pdf}
\caption{TrpB decision-cohort summaries: parameter resampling. Rows, top to bottom: trpb-libH, Gamma--Exponential; trpb-libI, Exponential--Exponential; trpb-libI, Lomax--Exponential. Columns show normalized central 90\% score width versus input count, output count, and median-score margin from the top-100 cutoff; then top-100 inclusion versus that margin. Widths and margins use the run's top-1,000 median-score interquartile range (IQR). Blue filled points are reference top-100 members; orange open points are outside candidates with inclusion probability at least 0.01. Count and score-width limits are shared within family and scheme; margin limits are shared within a row. Dotted lines mark the cutoff. No regression lines are fitted.}
\label{fig:why-parameter_posterior-trpb-9}
\end{figure}
\clearpage

\begin{figure}[p]
\centering
\includegraphics[width=\linewidth]{summary_why_parameter_posterior_trpb_10.pdf}
\caption{TrpB decision-cohort summaries: parameter resampling. Rows, top to bottom: trpb-libI, Gamma--Exponential. Columns show normalized central 90\% score width versus input count, output count, and median-score margin from the top-100 cutoff; then top-100 inclusion versus that margin. Widths and margins use the run's top-1,000 median-score interquartile range (IQR). Blue filled points are reference top-100 members; orange open points are outside candidates with inclusion probability at least 0.01. Count and score-width limits are shared within family and scheme; margin limits are shared within a row. Dotted lines mark the cutoff. No regression lines are fitted.}
\label{fig:why-parameter_posterior-trpb-10}
\end{figure}
\clearpage

\begin{figure}[p]
\centering
\includegraphics[width=\linewidth]{summary_why_parameter_posterior_nuc_01.pdf}
\caption{NucB decision-cohort summaries: parameter resampling. Rows, top to bottom: nuc-G1, Exponential; nuc-G1, Lomax; nuc-G2, Exponential. Columns show normalized central 90\% score width versus input count, output count, and median-score margin from the top-100 cutoff; then top-100 inclusion versus that margin. Widths and margins use the run's top-1,000 median-score interquartile range (IQR). Blue filled points are reference top-100 members; orange open points are outside candidates with inclusion probability at least 0.01. Count and score-width limits are shared within family and scheme; margin limits are shared within a row. Dotted lines mark the cutoff. No regression lines are fitted.}
\label{fig:why-parameter_posterior-nuc-1}
\end{figure}
\clearpage

\begin{figure}[p]
\centering
\includegraphics[width=\linewidth]{summary_why_parameter_posterior_nuc_02.pdf}
\caption{NucB decision-cohort summaries: parameter resampling. Rows, top to bottom: nuc-G2, Lomax; nuc-G3, Exponential; nuc-G4, Exponential. Columns show normalized central 90\% score width versus input count, output count, and median-score margin from the top-100 cutoff; then top-100 inclusion versus that margin. Widths and margins use the run's top-1,000 median-score interquartile range (IQR). Blue filled points are reference top-100 members; orange open points are outside candidates with inclusion probability at least 0.01. Count and score-width limits are shared within family and scheme; margin limits are shared within a row. Dotted lines mark the cutoff. No regression lines are fitted.}
\label{fig:why-parameter_posterior-nuc-2}
\end{figure}
\clearpage

\begin{figure}[p]
\centering
\includegraphics[width=\linewidth]{summary_why_parameter_posterior_nuc_03.pdf}
\caption{NucB decision-cohort summaries: parameter resampling. Rows, top to bottom: nuc-G4, Lomax. Columns show normalized central 90\% score width versus input count, output count, and median-score margin from the top-100 cutoff; then top-100 inclusion versus that margin. Widths and margins use the run's top-1,000 median-score interquartile range (IQR). Blue filled points are reference top-100 members; orange open points are outside candidates with inclusion probability at least 0.01. Count and score-width limits are shared within family and scheme; margin limits are shared within a row. Dotted lines mark the cutoff. No regression lines are fitted.}
\label{fig:why-parameter_posterior-nuc-3}
\end{figure}
\clearpage

\begin{figure}[p]
\centering
\includegraphics[width=\linewidth]{summary_why_mle_predictive_pdz_01.pdf}
\caption{PDZ decision-cohort summaries: optimized-fit count resampling. Rows, top to bottom: pdz-cis, Exponential--Exponential; pdz-cis, Lomax--Exponential; pdz-cis, Exponential--Sigmoid. Columns show normalized central 90\% score width versus input count, output count, and median-score margin from the top-100 cutoff; then top-100 inclusion versus that margin. Widths and margins use the run's top-1,000 median-score interquartile range (IQR). Blue filled points are reference top-100 members; orange open points are outside candidates with inclusion probability at least 0.01. Count and score-width limits are shared within family and scheme; margin limits are shared within a row. Dotted lines mark the cutoff. No regression lines are fitted.}
\label{fig:why-mle_predictive-pdz-1}
\end{figure}
\clearpage

\begin{figure}[p]
\centering
\includegraphics[width=\linewidth]{summary_why_mle_predictive_pdz_02.pdf}
\caption{PDZ decision-cohort summaries: optimized-fit count resampling. Rows, top to bottom: pdz-cript-c, Exponential--Exponential; pdz-cript-c, Lomax--Exponential; pdz-cript-c, Gamma--Exponential. Columns show normalized central 90\% score width versus input count, output count, and median-score margin from the top-100 cutoff; then top-100 inclusion versus that margin. Widths and margins use the run's top-1,000 median-score interquartile range (IQR). Blue filled points are reference top-100 members; orange open points are outside candidates with inclusion probability at least 0.01. Count and score-width limits are shared within family and scheme; margin limits are shared within a row. Dotted lines mark the cutoff. No regression lines are fitted.}
\label{fig:why-mle_predictive-pdz-2}
\end{figure}
\clearpage

\begin{figure}[p]
\centering
\includegraphics[width=\linewidth]{summary_why_mle_predictive_pdz_03.pdf}
\caption{PDZ decision-cohort summaries: optimized-fit count resampling. Rows, top to bottom: pdz-cript-c, Exponential--Sigmoid; pdz-cript-c, Lomax--Sigmoid; pdz-cript-n, Exponential--Exponential. Columns show normalized central 90\% score width versus input count, output count, and median-score margin from the top-100 cutoff; then top-100 inclusion versus that margin. Widths and margins use the run's top-1,000 median-score interquartile range (IQR). Blue filled points are reference top-100 members; orange open points are outside candidates with inclusion probability at least 0.01. Count and score-width limits are shared within family and scheme; margin limits are shared within a row. Dotted lines mark the cutoff. No regression lines are fitted.}
\label{fig:why-mle_predictive-pdz-3}
\end{figure}
\clearpage

\begin{figure}[p]
\centering
\includegraphics[width=\linewidth]{summary_why_mle_predictive_pdz_04.pdf}
\caption{PDZ decision-cohort summaries: optimized-fit count resampling. Rows, top to bottom: pdz-cript-n, Lomax--Exponential; pdz-cript-n, Gamma--Exponential; pdz-cript-n, Lomax--Sigmoid. Columns show normalized central 90\% score width versus input count, output count, and median-score margin from the top-100 cutoff; then top-100 inclusion versus that margin. Widths and margins use the run's top-1,000 median-score interquartile range (IQR). Blue filled points are reference top-100 members; orange open points are outside candidates with inclusion probability at least 0.01. Count and score-width limits are shared within family and scheme; margin limits are shared within a row. Dotted lines mark the cutoff. No regression lines are fitted.}
\label{fig:why-mle_predictive-pdz-4}
\end{figure}
\clearpage

\begin{figure}[p]
\centering
\includegraphics[width=\linewidth]{summary_why_mle_predictive_pdz_05.pdf}
\caption{PDZ decision-cohort summaries: optimized-fit count resampling. Rows, top to bottom: pdz-cript-n, Gamma--Sigmoid; pdz-trans-1, Exponential--Exponential; pdz-trans-1, Gamma--Exponential. Columns show normalized central 90\% score width versus input count, output count, and median-score margin from the top-100 cutoff; then top-100 inclusion versus that margin. Widths and margins use the run's top-1,000 median-score interquartile range (IQR). Blue filled points are reference top-100 members; orange open points are outside candidates with inclusion probability at least 0.01. Count and score-width limits are shared within family and scheme; margin limits are shared within a row. Dotted lines mark the cutoff. No regression lines are fitted.}
\label{fig:why-mle_predictive-pdz-5}
\end{figure}
\clearpage

\begin{figure}[p]
\centering
\includegraphics[width=\linewidth]{summary_why_mle_predictive_pdz_06.pdf}
\caption{PDZ decision-cohort summaries: optimized-fit count resampling. Rows, top to bottom: pdz-trans-2, Exponential--Exponential; pdz-trans-2, Gamma--Exponential; pdz-trans-2, Lomax--Sigmoid. Columns show normalized central 90\% score width versus input count, output count, and median-score margin from the top-100 cutoff; then top-100 inclusion versus that margin. Widths and margins use the run's top-1,000 median-score interquartile range (IQR). Blue filled points are reference top-100 members; orange open points are outside candidates with inclusion probability at least 0.01. Count and score-width limits are shared within family and scheme; margin limits are shared within a row. Dotted lines mark the cutoff. No regression lines are fitted.}
\label{fig:why-mle_predictive-pdz-6}
\end{figure}
\clearpage

\begin{figure}[p]
\centering
\includegraphics[width=\linewidth]{summary_why_mle_predictive_pdz_07.pdf}
\caption{PDZ decision-cohort summaries: optimized-fit count resampling. Rows, top to bottom: pdz-trans-2, Gamma--Sigmoid. Columns show normalized central 90\% score width versus input count, output count, and median-score margin from the top-100 cutoff; then top-100 inclusion versus that margin. Widths and margins use the run's top-1,000 median-score interquartile range (IQR). Blue filled points are reference top-100 members; orange open points are outside candidates with inclusion probability at least 0.01. Count and score-width limits are shared within family and scheme; margin limits are shared within a row. Dotted lines mark the cutoff. No regression lines are fitted.}
\label{fig:why-mle_predictive-pdz-7}
\end{figure}
\clearpage

\begin{figure}[p]
\centering
\includegraphics[width=\linewidth]{summary_why_mle_predictive_trpb_01.pdf}
\caption{TrpB decision-cohort summaries: optimized-fit count resampling. Rows, top to bottom: trpb-libA, Exponential--Exponential; trpb-libA, Lomax--Exponential; trpb-libA, Exponential--Sigmoid. Columns show normalized central 90\% score width versus input count, output count, and median-score margin from the top-100 cutoff; then top-100 inclusion versus that margin. Widths and margins use the run's top-1,000 median-score interquartile range (IQR). Blue filled points are reference top-100 members; orange open points are outside candidates with inclusion probability at least 0.01. Count and score-width limits are shared within family and scheme; margin limits are shared within a row. Dotted lines mark the cutoff. No regression lines are fitted.}
\label{fig:why-mle_predictive-trpb-1}
\end{figure}
\clearpage

\begin{figure}[p]
\centering
\includegraphics[width=\linewidth]{summary_why_mle_predictive_trpb_02.pdf}
\caption{TrpB decision-cohort summaries: optimized-fit count resampling. Rows, top to bottom: trpb-libB, Exponential--Exponential; trpb-libB, Lomax--Exponential; trpb-libB, Exponential--Sigmoid. Columns show normalized central 90\% score width versus input count, output count, and median-score margin from the top-100 cutoff; then top-100 inclusion versus that margin. Widths and margins use the run's top-1,000 median-score interquartile range (IQR). Blue filled points are reference top-100 members; orange open points are outside candidates with inclusion probability at least 0.01. Count and score-width limits are shared within family and scheme; margin limits are shared within a row. Dotted lines mark the cutoff. No regression lines are fitted.}
\label{fig:why-mle_predictive-trpb-2}
\end{figure}
\clearpage

\begin{figure}[p]
\centering
\includegraphics[width=\linewidth]{summary_why_mle_predictive_trpb_03.pdf}
\caption{TrpB decision-cohort summaries: optimized-fit count resampling. Rows, top to bottom: trpb-libC, Exponential--Exponential; trpb-libC, Lomax--Exponential; trpb-libC, Exponential--Sigmoid. Columns show normalized central 90\% score width versus input count, output count, and median-score margin from the top-100 cutoff; then top-100 inclusion versus that margin. Widths and margins use the run's top-1,000 median-score interquartile range (IQR). Blue filled points are reference top-100 members; orange open points are outside candidates with inclusion probability at least 0.01. Count and score-width limits are shared within family and scheme; margin limits are shared within a row. Dotted lines mark the cutoff. No regression lines are fitted.}
\label{fig:why-mle_predictive-trpb-3}
\end{figure}
\clearpage

\begin{figure}[p]
\centering
\includegraphics[width=\linewidth]{summary_why_mle_predictive_trpb_04.pdf}
\caption{TrpB decision-cohort summaries: optimized-fit count resampling. Rows, top to bottom: trpb-libD, Exponential--Exponential; trpb-libD, Lomax--Exponential; trpb-libD, Gamma--Exponential. Columns show normalized central 90\% score width versus input count, output count, and median-score margin from the top-100 cutoff; then top-100 inclusion versus that margin. Widths and margins use the run's top-1,000 median-score interquartile range (IQR). Blue filled points are reference top-100 members; orange open points are outside candidates with inclusion probability at least 0.01. Count and score-width limits are shared within family and scheme; margin limits are shared within a row. Dotted lines mark the cutoff. No regression lines are fitted.}
\label{fig:why-mle_predictive-trpb-4}
\end{figure}
\clearpage

\begin{figure}[p]
\centering
\includegraphics[width=\linewidth]{summary_why_mle_predictive_trpb_05.pdf}
\caption{TrpB decision-cohort summaries: optimized-fit count resampling. Rows, top to bottom: trpb-libE, Exponential--Exponential; trpb-libE, Lomax--Exponential; trpb-libE, Gamma--Exponential. Columns show normalized central 90\% score width versus input count, output count, and median-score margin from the top-100 cutoff; then top-100 inclusion versus that margin. Widths and margins use the run's top-1,000 median-score interquartile range (IQR). Blue filled points are reference top-100 members; orange open points are outside candidates with inclusion probability at least 0.01. Count and score-width limits are shared within family and scheme; margin limits are shared within a row. Dotted lines mark the cutoff. No regression lines are fitted.}
\label{fig:why-mle_predictive-trpb-5}
\end{figure}
\clearpage

\begin{figure}[p]
\centering
\includegraphics[width=\linewidth]{summary_why_mle_predictive_trpb_06.pdf}
\caption{TrpB decision-cohort summaries: optimized-fit count resampling. Rows, top to bottom: trpb-libF, Exponential--Exponential; trpb-libF, Lomax--Exponential; trpb-libF, Gamma--Exponential. Columns show normalized central 90\% score width versus input count, output count, and median-score margin from the top-100 cutoff; then top-100 inclusion versus that margin. Widths and margins use the run's top-1,000 median-score interquartile range (IQR). Blue filled points are reference top-100 members; orange open points are outside candidates with inclusion probability at least 0.01. Count and score-width limits are shared within family and scheme; margin limits are shared within a row. Dotted lines mark the cutoff. No regression lines are fitted.}
\label{fig:why-mle_predictive-trpb-6}
\end{figure}
\clearpage

\begin{figure}[p]
\centering
\includegraphics[width=\linewidth]{summary_why_mle_predictive_trpb_07.pdf}
\caption{TrpB decision-cohort summaries: optimized-fit count resampling. Rows, top to bottom: trpb-libF, Exponential--Sigmoid; trpb-libG, Exponential--Exponential; trpb-libG, Lomax--Exponential. Columns show normalized central 90\% score width versus input count, output count, and median-score margin from the top-100 cutoff; then top-100 inclusion versus that margin. Widths and margins use the run's top-1,000 median-score interquartile range (IQR). Blue filled points are reference top-100 members; orange open points are outside candidates with inclusion probability at least 0.01. Count and score-width limits are shared within family and scheme; margin limits are shared within a row. Dotted lines mark the cutoff. No regression lines are fitted.}
\label{fig:why-mle_predictive-trpb-7}
\end{figure}
\clearpage

\begin{figure}[p]
\centering
\includegraphics[width=\linewidth]{summary_why_mle_predictive_trpb_08.pdf}
\caption{TrpB decision-cohort summaries: optimized-fit count resampling. Rows, top to bottom: trpb-libG, Gamma--Exponential; trpb-libH, Exponential--Exponential; trpb-libH, Lomax--Exponential. Columns show normalized central 90\% score width versus input count, output count, and median-score margin from the top-100 cutoff; then top-100 inclusion versus that margin. Widths and margins use the run's top-1,000 median-score interquartile range (IQR). Blue filled points are reference top-100 members; orange open points are outside candidates with inclusion probability at least 0.01. Count and score-width limits are shared within family and scheme; margin limits are shared within a row. Dotted lines mark the cutoff. No regression lines are fitted.}
\label{fig:why-mle_predictive-trpb-8}
\end{figure}
\clearpage

\begin{figure}[p]
\centering
\includegraphics[width=\linewidth]{summary_why_mle_predictive_trpb_09.pdf}
\caption{TrpB decision-cohort summaries: optimized-fit count resampling. Rows, top to bottom: trpb-libH, Gamma--Exponential; trpb-libI, Exponential--Exponential; trpb-libI, Lomax--Exponential. Columns show normalized central 90\% score width versus input count, output count, and median-score margin from the top-100 cutoff; then top-100 inclusion versus that margin. Widths and margins use the run's top-1,000 median-score interquartile range (IQR). Blue filled points are reference top-100 members; orange open points are outside candidates with inclusion probability at least 0.01. Count and score-width limits are shared within family and scheme; margin limits are shared within a row. Dotted lines mark the cutoff. No regression lines are fitted.}
\label{fig:why-mle_predictive-trpb-9}
\end{figure}
\clearpage

\begin{figure}[p]
\centering
\includegraphics[width=\linewidth]{summary_why_mle_predictive_trpb_10.pdf}
\caption{TrpB decision-cohort summaries: optimized-fit count resampling. Rows, top to bottom: trpb-libI, Gamma--Exponential; trpb-four-site, Exponential--Exponential; trpb-four-site, Lomax--Exponential. Columns show normalized central 90\% score width versus input count, output count, and median-score margin from the top-100 cutoff; then top-100 inclusion versus that margin. Widths and margins use the run's top-1,000 median-score interquartile range (IQR). Blue filled points are reference top-100 members; orange open points are outside candidates with inclusion probability at least 0.01. Count and score-width limits are shared within family and scheme; margin limits are shared within a row. Dotted lines mark the cutoff. No regression lines are fitted.}
\label{fig:why-mle_predictive-trpb-10}
\end{figure}
\clearpage

\begin{figure}[p]
\centering
\includegraphics[width=\linewidth]{summary_why_mle_predictive_trpb_11.pdf}
\caption{TrpB decision-cohort summaries: optimized-fit count resampling. Rows, top to bottom: trpb-four-site, Gamma--Exponential; trpb-four-site, Exponential--Sigmoid; trpb-four-site, Lomax--Sigmoid. Columns show normalized central 90\% score width versus input count, output count, and median-score margin from the top-100 cutoff; then top-100 inclusion versus that margin. Widths and margins use the run's top-1,000 median-score interquartile range (IQR). Blue filled points are reference top-100 members; orange open points are outside candidates with inclusion probability at least 0.01. Count and score-width limits are shared within family and scheme; margin limits are shared within a row. Dotted lines mark the cutoff. No regression lines are fitted.}
\label{fig:why-mle_predictive-trpb-11}
\end{figure}
\clearpage

\begin{figure}[p]
\centering
\includegraphics[width=\linewidth]{summary_why_mle_predictive_trpb_12.pdf}
\caption{TrpB decision-cohort summaries: optimized-fit count resampling. Rows, top to bottom: trpb-four-site, Gamma--Sigmoid. Columns show normalized central 90\% score width versus input count, output count, and median-score margin from the top-100 cutoff; then top-100 inclusion versus that margin. Widths and margins use the run's top-1,000 median-score interquartile range (IQR). Blue filled points are reference top-100 members; orange open points are outside candidates with inclusion probability at least 0.01. Count and score-width limits are shared within family and scheme; margin limits are shared within a row. Dotted lines mark the cutoff. No regression lines are fitted.}
\label{fig:why-mle_predictive-trpb-12}
\end{figure}
\clearpage

\begin{figure}[p]
\centering
\includegraphics[width=\linewidth]{summary_why_mle_predictive_nuc_01.pdf}
\caption{NucB decision-cohort summaries: optimized-fit count resampling. Rows, top to bottom: nuc-G1, Exponential; nuc-G1, Lomax; nuc-G2, Exponential. Columns show normalized central 90\% score width versus input count, output count, and median-score margin from the top-100 cutoff; then top-100 inclusion versus that margin. Widths and margins use the run's top-1,000 median-score interquartile range (IQR). Blue filled points are reference top-100 members; orange open points are outside candidates with inclusion probability at least 0.01. Count and score-width limits are shared within family and scheme; margin limits are shared within a row. Dotted lines mark the cutoff. No regression lines are fitted.}
\label{fig:why-mle_predictive-nuc-1}
\end{figure}
\clearpage

\begin{figure}[p]
\centering
\includegraphics[width=\linewidth]{summary_why_mle_predictive_nuc_02.pdf}
\caption{NucB decision-cohort summaries: optimized-fit count resampling. Rows, top to bottom: nuc-G2, Lomax; nuc-G3, Exponential; nuc-G4, Exponential. Columns show normalized central 90\% score width versus input count, output count, and median-score margin from the top-100 cutoff; then top-100 inclusion versus that margin. Widths and margins use the run's top-1,000 median-score interquartile range (IQR). Blue filled points are reference top-100 members; orange open points are outside candidates with inclusion probability at least 0.01. Count and score-width limits are shared within family and scheme; margin limits are shared within a row. Dotted lines mark the cutoff. No regression lines are fitted.}
\label{fig:why-mle_predictive-nuc-2}
\end{figure}
\clearpage

\begin{figure}[p]
\centering
\includegraphics[width=\linewidth]{summary_why_mle_predictive_nuc_03.pdf}
\caption{NucB decision-cohort summaries: optimized-fit count resampling. Rows, top to bottom: nuc-G4, Lomax. Columns show normalized central 90\% score width versus input count, output count, and median-score margin from the top-100 cutoff; then top-100 inclusion versus that margin. Widths and margins use the run's top-1,000 median-score interquartile range (IQR). Blue filled points are reference top-100 members; orange open points are outside candidates with inclusion probability at least 0.01. Count and score-width limits are shared within family and scheme; margin limits are shared within a row. Dotted lines mark the cutoff. No regression lines are fitted.}
\label{fig:why-mle_predictive-nuc-3}
\end{figure}
\clearpage

\begin{figure}[p]
\centering
\includegraphics[width=\linewidth]{summary_why_posterior_predictive_pdz_01.pdf}
\caption{PDZ decision-cohort summaries: posterior count resampling. Rows, top to bottom: pdz-trans-1, Exponential--Exponential; pdz-trans-1, Gamma--Exponential; pdz-trans-2, Exponential--Exponential. Columns show normalized central 90\% score width versus input count, output count, and median-score margin from the top-100 cutoff; then top-100 inclusion versus that margin. Widths and margins use the run's top-1,000 median-score interquartile range (IQR). Blue filled points are reference top-100 members; orange open points are outside candidates with inclusion probability at least 0.01. Count and score-width limits are shared within family and scheme; margin limits are shared within a row. Dotted lines mark the cutoff. No regression lines are fitted.}
\label{fig:why-posterior_predictive-pdz-1}
\end{figure}
\clearpage

\begin{figure}[p]
\centering
\includegraphics[width=\linewidth]{summary_why_posterior_predictive_pdz_02.pdf}
\caption{PDZ decision-cohort summaries: posterior count resampling. Rows, top to bottom: pdz-trans-2, Gamma--Sigmoid; pdz-cis, Exponential--Exponential; pdz-cis, Lomax--Exponential. Columns show normalized central 90\% score width versus input count, output count, and median-score margin from the top-100 cutoff; then top-100 inclusion versus that margin. Widths and margins use the run's top-1,000 median-score interquartile range (IQR). Blue filled points are reference top-100 members; orange open points are outside candidates with inclusion probability at least 0.01. Count and score-width limits are shared within family and scheme; margin limits are shared within a row. Dotted lines mark the cutoff. No regression lines are fitted.}
\label{fig:why-posterior_predictive-pdz-2}
\end{figure}
\clearpage

\begin{figure}[p]
\centering
\includegraphics[width=\linewidth]{summary_why_posterior_predictive_pdz_03.pdf}
\caption{PDZ decision-cohort summaries: posterior count resampling. Rows, top to bottom: pdz-cis, Exponential--Sigmoid; pdz-cript-c, Exponential--Exponential; pdz-cript-c, Lomax--Exponential. Columns show normalized central 90\% score width versus input count, output count, and median-score margin from the top-100 cutoff; then top-100 inclusion versus that margin. Widths and margins use the run's top-1,000 median-score interquartile range (IQR). Blue filled points are reference top-100 members; orange open points are outside candidates with inclusion probability at least 0.01. Count and score-width limits are shared within family and scheme; margin limits are shared within a row. Dotted lines mark the cutoff. No regression lines are fitted.}
\label{fig:why-posterior_predictive-pdz-3}
\end{figure}
\clearpage

\begin{figure}[p]
\centering
\includegraphics[width=\linewidth]{summary_why_posterior_predictive_pdz_04.pdf}
\caption{PDZ decision-cohort summaries: posterior count resampling. Rows, top to bottom: pdz-cript-c, Gamma--Exponential; pdz-cript-c, Exponential--Sigmoid; pdz-cript-n, Gamma--Exponential. Columns show normalized central 90\% score width versus input count, output count, and median-score margin from the top-100 cutoff; then top-100 inclusion versus that margin. Widths and margins use the run's top-1,000 median-score interquartile range (IQR). Blue filled points are reference top-100 members; orange open points are outside candidates with inclusion probability at least 0.01. Count and score-width limits are shared within family and scheme; margin limits are shared within a row. Dotted lines mark the cutoff. No regression lines are fitted.}
\label{fig:why-posterior_predictive-pdz-4}
\end{figure}
\clearpage

\begin{figure}[p]
\centering
\includegraphics[width=\linewidth]{summary_why_posterior_predictive_pdz_05.pdf}
\caption{PDZ decision-cohort summaries: posterior count resampling. Rows, top to bottom: pdz-cript-n, Lomax--Sigmoid; pdz-cript-n, Gamma--Sigmoid. Columns show normalized central 90\% score width versus input count, output count, and median-score margin from the top-100 cutoff; then top-100 inclusion versus that margin. Widths and margins use the run's top-1,000 median-score interquartile range (IQR). Blue filled points are reference top-100 members; orange open points are outside candidates with inclusion probability at least 0.01. Count and score-width limits are shared within family and scheme; margin limits are shared within a row. Dotted lines mark the cutoff. No regression lines are fitted.}
\label{fig:why-posterior_predictive-pdz-5}
\end{figure}
\clearpage

\begin{figure}[p]
\centering
\includegraphics[width=\linewidth]{summary_why_posterior_predictive_trpb_01.pdf}
\caption{TrpB decision-cohort summaries: posterior count resampling. Rows, top to bottom: trpb-libA, Exponential--Exponential; trpb-libA, Lomax--Exponential; trpb-libA, Exponential--Sigmoid. Columns show normalized central 90\% score width versus input count, output count, and median-score margin from the top-100 cutoff; then top-100 inclusion versus that margin. Widths and margins use the run's top-1,000 median-score interquartile range (IQR). Blue filled points are reference top-100 members; orange open points are outside candidates with inclusion probability at least 0.01. Count and score-width limits are shared within family and scheme; margin limits are shared within a row. Dotted lines mark the cutoff. No regression lines are fitted.}
\label{fig:why-posterior_predictive-trpb-1}
\end{figure}
\clearpage

\begin{figure}[p]
\centering
\includegraphics[width=\linewidth]{summary_why_posterior_predictive_trpb_02.pdf}
\caption{TrpB decision-cohort summaries: posterior count resampling. Rows, top to bottom: trpb-libB, Exponential--Exponential; trpb-libB, Lomax--Exponential; trpb-libB, Exponential--Sigmoid. Columns show normalized central 90\% score width versus input count, output count, and median-score margin from the top-100 cutoff; then top-100 inclusion versus that margin. Widths and margins use the run's top-1,000 median-score interquartile range (IQR). Blue filled points are reference top-100 members; orange open points are outside candidates with inclusion probability at least 0.01. Count and score-width limits are shared within family and scheme; margin limits are shared within a row. Dotted lines mark the cutoff. No regression lines are fitted.}
\label{fig:why-posterior_predictive-trpb-2}
\end{figure}
\clearpage

\begin{figure}[p]
\centering
\includegraphics[width=\linewidth]{summary_why_posterior_predictive_trpb_03.pdf}
\caption{TrpB decision-cohort summaries: posterior count resampling. Rows, top to bottom: trpb-libC, Exponential--Exponential; trpb-libC, Lomax--Exponential; trpb-libC, Exponential--Sigmoid. Columns show normalized central 90\% score width versus input count, output count, and median-score margin from the top-100 cutoff; then top-100 inclusion versus that margin. Widths and margins use the run's top-1,000 median-score interquartile range (IQR). Blue filled points are reference top-100 members; orange open points are outside candidates with inclusion probability at least 0.01. Count and score-width limits are shared within family and scheme; margin limits are shared within a row. Dotted lines mark the cutoff. No regression lines are fitted.}
\label{fig:why-posterior_predictive-trpb-3}
\end{figure}
\clearpage

\begin{figure}[p]
\centering
\includegraphics[width=\linewidth]{summary_why_posterior_predictive_trpb_04.pdf}
\caption{TrpB decision-cohort summaries: posterior count resampling. Rows, top to bottom: trpb-libD, Exponential--Exponential; trpb-libD, Lomax--Exponential; trpb-libD, Gamma--Exponential. Columns show normalized central 90\% score width versus input count, output count, and median-score margin from the top-100 cutoff; then top-100 inclusion versus that margin. Widths and margins use the run's top-1,000 median-score interquartile range (IQR). Blue filled points are reference top-100 members; orange open points are outside candidates with inclusion probability at least 0.01. Count and score-width limits are shared within family and scheme; margin limits are shared within a row. Dotted lines mark the cutoff. No regression lines are fitted.}
\label{fig:why-posterior_predictive-trpb-4}
\end{figure}
\clearpage

\begin{figure}[p]
\centering
\includegraphics[width=\linewidth]{summary_why_posterior_predictive_trpb_05.pdf}
\caption{TrpB decision-cohort summaries: posterior count resampling. Rows, top to bottom: trpb-libE, Exponential--Exponential; trpb-libE, Lomax--Exponential; trpb-libE, Gamma--Exponential. Columns show normalized central 90\% score width versus input count, output count, and median-score margin from the top-100 cutoff; then top-100 inclusion versus that margin. Widths and margins use the run's top-1,000 median-score interquartile range (IQR). Blue filled points are reference top-100 members; orange open points are outside candidates with inclusion probability at least 0.01. Count and score-width limits are shared within family and scheme; margin limits are shared within a row. Dotted lines mark the cutoff. No regression lines are fitted.}
\label{fig:why-posterior_predictive-trpb-5}
\end{figure}
\clearpage

\begin{figure}[p]
\centering
\includegraphics[width=\linewidth]{summary_why_posterior_predictive_trpb_06.pdf}
\caption{TrpB decision-cohort summaries: posterior count resampling. Rows, top to bottom: trpb-libF, Exponential--Exponential; trpb-libF, Lomax--Exponential; trpb-libF, Gamma--Exponential. Columns show normalized central 90\% score width versus input count, output count, and median-score margin from the top-100 cutoff; then top-100 inclusion versus that margin. Widths and margins use the run's top-1,000 median-score interquartile range (IQR). Blue filled points are reference top-100 members; orange open points are outside candidates with inclusion probability at least 0.01. Count and score-width limits are shared within family and scheme; margin limits are shared within a row. Dotted lines mark the cutoff. No regression lines are fitted.}
\label{fig:why-posterior_predictive-trpb-6}
\end{figure}
\clearpage

\begin{figure}[p]
\centering
\includegraphics[width=\linewidth]{summary_why_posterior_predictive_trpb_07.pdf}
\caption{TrpB decision-cohort summaries: posterior count resampling. Rows, top to bottom: trpb-libF, Exponential--Sigmoid; trpb-libG, Exponential--Exponential; trpb-libG, Lomax--Exponential. Columns show normalized central 90\% score width versus input count, output count, and median-score margin from the top-100 cutoff; then top-100 inclusion versus that margin. Widths and margins use the run's top-1,000 median-score interquartile range (IQR). Blue filled points are reference top-100 members; orange open points are outside candidates with inclusion probability at least 0.01. Count and score-width limits are shared within family and scheme; margin limits are shared within a row. Dotted lines mark the cutoff. No regression lines are fitted.}
\label{fig:why-posterior_predictive-trpb-7}
\end{figure}
\clearpage

\begin{figure}[p]
\centering
\includegraphics[width=\linewidth]{summary_why_posterior_predictive_trpb_08.pdf}
\caption{TrpB decision-cohort summaries: posterior count resampling. Rows, top to bottom: trpb-libG, Gamma--Exponential; trpb-libH, Exponential--Exponential; trpb-libH, Lomax--Exponential. Columns show normalized central 90\% score width versus input count, output count, and median-score margin from the top-100 cutoff; then top-100 inclusion versus that margin. Widths and margins use the run's top-1,000 median-score interquartile range (IQR). Blue filled points are reference top-100 members; orange open points are outside candidates with inclusion probability at least 0.01. Count and score-width limits are shared within family and scheme; margin limits are shared within a row. Dotted lines mark the cutoff. No regression lines are fitted.}
\label{fig:why-posterior_predictive-trpb-8}
\end{figure}
\clearpage

\begin{figure}[p]
\centering
\includegraphics[width=\linewidth]{summary_why_posterior_predictive_trpb_09.pdf}
\caption{TrpB decision-cohort summaries: posterior count resampling. Rows, top to bottom: trpb-libH, Gamma--Exponential; trpb-libI, Exponential--Exponential; trpb-libI, Lomax--Exponential. Columns show normalized central 90\% score width versus input count, output count, and median-score margin from the top-100 cutoff; then top-100 inclusion versus that margin. Widths and margins use the run's top-1,000 median-score interquartile range (IQR). Blue filled points are reference top-100 members; orange open points are outside candidates with inclusion probability at least 0.01. Count and score-width limits are shared within family and scheme; margin limits are shared within a row. Dotted lines mark the cutoff. No regression lines are fitted.}
\label{fig:why-posterior_predictive-trpb-9}
\end{figure}
\clearpage

\begin{figure}[p]
\centering
\includegraphics[width=\linewidth]{summary_why_posterior_predictive_trpb_10.pdf}
\caption{TrpB decision-cohort summaries: posterior count resampling. Rows, top to bottom: trpb-libI, Gamma--Exponential. Columns show normalized central 90\% score width versus input count, output count, and median-score margin from the top-100 cutoff; then top-100 inclusion versus that margin. Widths and margins use the run's top-1,000 median-score interquartile range (IQR). Blue filled points are reference top-100 members; orange open points are outside candidates with inclusion probability at least 0.01. Count and score-width limits are shared within family and scheme; margin limits are shared within a row. Dotted lines mark the cutoff. No regression lines are fitted.}
\label{fig:why-posterior_predictive-trpb-10}
\end{figure}
\clearpage

\begin{figure}[p]
\centering
\includegraphics[width=\linewidth]{summary_why_posterior_predictive_nuc_01.pdf}
\caption{NucB decision-cohort summaries: posterior count resampling. Rows, top to bottom: nuc-G1, Exponential; nuc-G1, Lomax; nuc-G2, Exponential. Columns show normalized central 90\% score width versus input count, output count, and median-score margin from the top-100 cutoff; then top-100 inclusion versus that margin. Widths and margins use the run's top-1,000 median-score interquartile range (IQR). Blue filled points are reference top-100 members; orange open points are outside candidates with inclusion probability at least 0.01. Count and score-width limits are shared within family and scheme; margin limits are shared within a row. Dotted lines mark the cutoff. No regression lines are fitted.}
\label{fig:why-posterior_predictive-nuc-1}
\end{figure}
\clearpage

\begin{figure}[p]
\centering
\includegraphics[width=\linewidth]{summary_why_posterior_predictive_nuc_02.pdf}
\caption{NucB decision-cohort summaries: posterior count resampling. Rows, top to bottom: nuc-G2, Lomax; nuc-G3, Exponential; nuc-G4, Exponential. Columns show normalized central 90\% score width versus input count, output count, and median-score margin from the top-100 cutoff; then top-100 inclusion versus that margin. Widths and margins use the run's top-1,000 median-score interquartile range (IQR). Blue filled points are reference top-100 members; orange open points are outside candidates with inclusion probability at least 0.01. Count and score-width limits are shared within family and scheme; margin limits are shared within a row. Dotted lines mark the cutoff. No regression lines are fitted.}
\label{fig:why-posterior_predictive-nuc-2}
\end{figure}
\clearpage

\begin{figure}[p]
\centering
\includegraphics[width=\linewidth]{summary_why_posterior_predictive_nuc_03.pdf}
\caption{NucB decision-cohort summaries: posterior count resampling. Rows, top to bottom: nuc-G4, Lomax. Columns show normalized central 90\% score width versus input count, output count, and median-score margin from the top-100 cutoff; then top-100 inclusion versus that margin. Widths and margins use the run's top-1,000 median-score interquartile range (IQR). Blue filled points are reference top-100 members; orange open points are outside candidates with inclusion probability at least 0.01. Count and score-width limits are shared within family and scheme; margin limits are shared within a row. Dotted lines mark the cutoff. No regression lines are fitted.}
\label{fig:why-posterior_predictive-nuc-3}
\end{figure}
\clearpage
